% concatenated from IYB-group-class-two.tex
\documentclass[12pt,a4paper]{article}
%\documentclass[12pt, reqno]{amsart}
\usepackage[T1]{fontenc}
\usepackage[british]{babel}
\usepackage[utf8]{inputenc}
\usepackage{amsmath,amssymb,amsthm}
%\usepackage{mathabx}
%\usepackage{paralist}
\usepackage{hyperref}

\newcommand\Aa{\mathbf{A}}
\newcommand\Bb{\mathbf{B}}
\newcommand\Ff{\mathfrak{F}}
\newcommand\FF{\mathbb{F}}
\newcommand\ZZ{\mathbf{Z}}
\newcommand\OO{\mathcal{O}}
\newcommand\GG{\mathfrak{G}}
\newcommand\HH{\mathfrak{H}}
\newcommand\Ss{\mathfrak{S}}
\newcommand\Ee{\mathfrak{E}}
\newcommand\XX{\mathfrak{X}}
\newcommand\F{\mathrm{F}}
\newcommand\K{\mathrm{K}}
\newcommand\PP{\mathbb{P}}
\newcommand\NN{\mathfrak{N}}
\newcommand{\clD}{\operatorname{\textsc{d}}}
\newcommand{\clS}{\operatorname{\textsc{s}}}

\DeclareMathOperator{\Crit}{Crit}
\DeclareMathOperator{\sn}{sn}
\DeclareMathOperator{\N}{N}
\DeclareMathOperator{\Oo}{O}
\DeclareMathOperator{\Hall}{Hall}
\DeclareMathOperator{\C}{C}
\DeclareMathOperator{\Soc}{Soc}
\DeclareMathOperator{\Core}{Core}
\DeclareMathOperator{\Z}{Z}
\DeclareMathOperator{\Fit}{F}
\DeclareMathOperator{\Ker}{Ker}
\DeclareMathOperator{\Exp}{exp}
\DeclareMathOperator{\Char}{Char}
\DeclareMathOperator{\Syl}{Syl}
\DeclareMathOperator{\Sylow}{Sylow}
\DeclareMathOperator{\Aut}{Aut}
\DeclareMathOperator{\Id}{Id}
\DeclareMathOperator{\GF}{GF}
\DeclareMathOperator{\Stab}{Stab}
\DeclareMathOperator{\Q}{Q}
\DeclareMathOperator{\Dim}{dim}
\DeclareMathOperator{\GL}{GL}
\DeclareMathOperator{\SL}{SL}
\DeclareMathOperator{\IYB}{IYB}
\DeclareMathOperator{\YB}{YB}
\DeclareMathOperator{\PSL}{PSL}
\DeclareMathOperator{\PGL}{PGL}
\DeclareMathOperator{\End}{End}
\DeclareMathOperator{\Mat}{Mat}
\DeclareMathOperator{\Gal}{Gal}
\DeclareMathOperator{\Out}{Out}
\DeclareMathOperator{\Mov}{Mov}
\DeclareMathOperator{\Sym}{Sym}
\DeclareMathOperator{\AGL}{AGL}
\newtheorem{theorem}{Theorem}
\newtheorem{corollary}[theorem]{Corollary}
\newtheorem{lemma}[theorem]{Lemma}
\newtheorem{question}[theorem]{Question}
\newtheorem{athm}{Theorem}
\renewcommand{\theathm}{\Alph{athm}}
\theoremstyle{definition}
\newtheorem{definition}[theorem]{Definition}
%\newtheorem{corollary}[theorem]{Corollary}
%\newtheorem{question}[theorem]{Question}
\newtheorem{problem}[theorem]{Problem}
\newtheorem{example}[theorem]{Example}

\title{On a class of involutive Yang-Baxter groups}
%\small{}}

\author{H. Meng\thanks{E-mail: hymeng2009@shu.edu.cn. Department of Mathematics
and Newtouch Center for Mathematics of Shanghai University,
Shanghai 200444, P. R. China. },~A. Ballester-Bolinches\thanks{E-mail: Adolfo.Ballester@uv.es. Departament de Matem\'{a}tiques, Universitat de Val\'{e}ncia, Dr. Moliner 50, 46100 Burjassot, Val\'{e}ncia,
Spain.},~P. Jim\'{e}nez-Seral\thanks{E-mail: paz@unizar.es. Departamento de Matem\'{a}ticas, Universidad de Zaragoza, Pedro Cerbuna, 12, 50009 Zaragoza, Spain.}
}
\date{}
\begin{document}
\maketitle

\begin{abstract}
A group is called an involutive Yang-Baxter group (IYB-group) if it is isomorphic to the permutation group of an involutive, non-degenerate set-theoretic solution of the Yang-Baxter equation. This paper investigates finite soluble groups whose Sylow subgroups have nilpotency class at most two, addressing Ced\'{o} and Okni\'{n}ski's question~\cite{CedoOkninski2025} of whether such groups are IYB-groups. We establish that a finite soluble group with Sylow subgroups of class at most two is an IYB-group if its nilpotent residual is $\Q_8$-free. We also prove that a finite soluble group with Sylow subgroups of class at most two and Sylow $2$-subgroups isomorphic to $Q_8$ is an IYB-group. Our approach depends on a careful study of the arithmetic and normal structure of a finite soluble group which is of independent interest.

\noindent\emph{Mathematics Subject Classification (2010): 16T25, 20D10, 20F16.}

\noindent\emph{Keywords: left braces, Yang–Baxter
equation, soluble groups, system normalisers}
\end{abstract}

\section{Introduction}
The Yang-Baxter equation plays a fundamental role in mathematical physics and serves as a cornerstone in the theory of quantum groups. In~\cite{Drinfeld1992}, Drinfeld proposed investigating set-theoretic solutions to this equation. A set-theoretic solution of the Yang–Baxter
equation is a pair $(X,r)$, where $X$ is a nonempty set and $r: X\times X \rightarrow X \times X$ is a map such that
$$r_{12}r_{23}r_{12}=r_{23}r_{12}r_{23},$$
where $r_{12}=r\times \Id_X, r_{23}=\Id_X \times r$ are maps from $X^3$ to $X^3$. Subsequent research by Etingof, Schedler and Soloviev~\cite{EtingofSchedlerSoloviev1999}, as well as Gateva-Ivanova and Van den Bergh~\cite{Gateva-IvanovaVandenBergh1998},
led to the study of a specific family of solutions known as \emph{involutive non-degenerate set-theoretic solutions.}

For all \( x, y \in X \), we define two maps \( f_x : X \longrightarrow X \) and \( g_y : X \longrightarrow X \) by setting \( r(x, y) = (f_x(y), g_y(x)) \). We say that the solution \((X, r)\) is \emph{involutive} if \( r^2 = \text{id}_{X \times X} \), and that \((X, r)\) is \emph{non-degenerate} if \( f_x, g_y \) are bijective maps for all \( x, y \in X \). By a \emph{solution of the YBE}, we will understand an involutive, non-degenerate set-theoretic solution of the Yang-Baxter equation.

Let \((X, r)\) be a solution of the YBE. The \emph{permutation group} of \((X, r)\) is the subgroup \(\mathcal{G}(X, r)\) of \(\text{Sym}(X)\) generated by the bijections \( f_x \) for all \( x \in X \), that is,
\[
\mathcal{G}(X, r) = \langle f_x \mid x \in X \rangle \leq \text{Sym}(X).
\]
A group \( G \) is called an \emph{involutive Yang-Baxter group}, or an IYB-group for short, if there exists a solution \((X, r)\) of the YBE such that \( G \cong \mathcal{G}(X, r) \).


To further explore such solutions, Rump~\cite{Rump2007} introduced a new algebraic structure called left brace, which was extended to a non-commutative setting by Guarnieri and Vendramin in \cite{GuarnieriVendramin2017}.

A \emph{skew (left) brace} is a triple \((B, +, \cdot)\) where \((B, +)\) and \((B, \cdot)\) are groups, referred to as the additive group (non-necessarily abelian) and the multiplicative group respectively, satisfying the following compatibility condition for all \(a, b, c \in B\):
\[
a \cdot (b + c) = a \cdot b - a + a \cdot c.
\]
Given a class of groups $\mathfrak{X}$, a skew left brace \(B\) is said to be of $\mathfrak{X}$-type if its additive group belongs to $\mathfrak{X}$. Rump's left braces are exactly the skew left braces of $\mathcal{A}$-type, where $\mathcal{A}$ is the class of all abelian groups. 


Following \cite{BEJP2023}, a group $G$ is called an $\mathfrak{X}\YB$-group if $G$ is isomorphic to the multiplicative group of a skew left brace of $\mathfrak{X}$-type. 

Applying \cite[Theorem~2.1]{CedoJespersRio2010}, it follows that a group $(G, \cdot)$ is an IYB-group if and only if it can be defined a binary operation $+$ on $G$ such that $(G,+,\cdot)$ is a left brace, that is, $(G, \cdot)$ is an $\mathcal{A}\YB$-group.


If $(G,\cdot)$ is a group in $\mathfrak{X}$, then we can regard $(G,+, \cdot)$ as a \emph{trivial skew left brace} of  $\mathfrak{X}$-type with with the addition "+" defined by $a+b:=ab, a,b \in G.$ Hence $\mathfrak{X}$ is a subclass of the class of all $\mathfrak{X}\YB$-groups.


 %A \emph{left brace} is defined as a set $B$ with two binary operations, $"+"$ and $"\cdot"$ such that $(B,+)$ is an abelian group, $(B,\cdot)$ is a group and
%$a\cdot (b+c)=a\cdot b+ a \cdot c-a$ holds for any $a,b,c \in B.$






The class of all finite IYB-groups is a subclass of the class $\mathcal S$ of all finite soluble groups (\cite{EtingofSchedlerSoloviev1999}) and, by a result of Bachiller (\cite{Bachiller2016}), it is a proper subclass of $\mathcal S$. However, it includes the class of all finite nilpotent groups of class two (\cite{AultWatters1973}), the class of all finite abelian-by-cyclic groups (\cite[Corollary~3.10]{CedoJespersRio2010}), and the class of finite soluble
groups which are semidirect products $A \rtimes H$, where $A$ is nilpotent of class at most two with odd order commutator subgroup and $H$ is an IYB-group (\cite[Corollary~4.2 and Corollary~4.3]{MengBallesterEstebanFuster2021} and \cite[Proposition~2.2]{CedoOkninski2025}]).


%The wreath product construction also preserves this property when both factors are IYB-groups~\cite[Theorem~3.5]{CedoJespersRio2010}, and finite

Another important class of IYB-groups is the class of all finite soluble groups whose Sylow subgroups are abelian (\cite[Corollary~4.3]{DavidGinosar2016}). Recently, Ced\'{o} and Okni\'{n}ski proved the following significant improvement of the above result.

%Research has established several conditions under which a finite group can be realised as the multiplicative group of a left brace. It is known that all abelian groups arise in this way, corresponding to trivial braces, and that this property is preserved under direct products. More generally, several important classes of groups have been shown to have this property: nilpotent groups of class two~\cite{AultWatters1973}, abelian-by-cyclic groups~\cite[Corollary~3.10]{CedoJespersRio2010}. In~\cite[Theorem~3.3]{CedoJespersRio2010}, semidirect products $A \rtimes H$ where $A$ is abelian and $H$ itself comes from a brace,
%moreover, the authors \cite[Corollary~4.2,4.3]{MengBallesterEstebanFuster2021} extend this result to the case where $A$ has nilpotency class two and possesses Abelian Sylow $2$-subgroups. The wreath product construction also preserves this property when both factors are IYB-groups~\cite[Theorem~3.5]{CedoJespersRio2010}. Furthermore, soluble groups whose Sylow subgroups are all abelian (so-called $A$-type groups) can also be IYB-groups~\cite[Corollary~4.3]{DavidGinosar2016}. ,


\begin{theorem}[{\cite[Theorem~2.6]{CedoOkninski2025}}]
Let $G$ be a finite soluble group with all Sylow subgroups of nilpotency class at most two, and Sylow $2$-subgroups abelian. Then $G$ is an $\IYB$-group.
\end{theorem}




A natural and interesting question~\cite[Question~2.7]{CedoOkninski2025} is also raised:
\begin{question}
Let $G$ be a finite soluble group such that all Sylow subgroups have nilpotency class at most two. Is $G$ an IYB-group?
\end{question}

We will address this question by establishing several significant results that contribute to the solution of this question. Our first result extends Ced\'{o} and Okni\'{n}ski's result to a wider class of all finite soluble groups with Sylow subgroups of class at most $2$.

It is well-known that the class $\mathcal{N}$ of all nilpotent groups is a class of groups which is closed under taking epimorphic images and finite subdirect products. Hence every finite group $G$ has a smallest normal subgroup with nilpotent quotient. This subgroup is called  the \emph{nilpotent residual} of $G$ and it is denoted by $G^{\mathcal{N}}$.


A group $G$ is said to be $\Q_8$-\emph{free} if $G$ has no section isomorphic to the quaternion group $Q_8$ of order $8$.


%: the case of quaternion-free groups ($\Q_8$-free groups for short)., or groups $G$ with no section $K/H$ ($H \unlhd K \leq G$) isomorphic to the quaternion group $\Q_8$ of order $8$.

\begin{athm}\label{thm-Q-8}
Let $G$ be a finite soluble group such that all Sylow subgroups have nilpotency class at most two. If  $G^{\mathcal{N}}$ is $\Q_8$-free, then $G$ is an $\IYB$-group.
\end{athm}


In~\cite[Theorem~2.3]{CedoOkninski2025},  Ced\'{o} and Okni\'{n}ski proved that $G$ is a finite soluble Sylow tower group with all Sylow subgroups of nilpotency class at most two, then $G$ is an IYB-group provided that $G$ is either of odd order or the Sylow $2$-subgroups of $G$ are abelian. As a consequence of Theorem~\ref{thm-Q-8}, a related result holds under the weaker assumption that $G$ is $2$-nilpotent as the nilpotent residual $G^{\mathcal{N}}$ of $G$ is of odd order.


\begin{corollary}
Let $G$ be a finite soluble group with all Sylow subgroups of nilpotency class at most two. Suppose that $G$ is $2$-nilpotent. Then $G$ is an $\IYB$-group.
\end{corollary}


If $G$ is a finite soluble group with Sylow subgroups of nilpotency class at most two and Sylow $2$-subgroups isomorphic to the dihedral group of order $8$, then $G$ is an IYB-group by Theorem~\ref{thm-Q-8}. Our next result analyses the $\Q_8$-case. The reader is referred to \cite[Theorem~3.1]{Bachiller2015-p3} for the classification of left braces of order~$8$.

\begin{athm}\label{thm-Sylow-Q-8}
Let $G$ be a finite soluble group with all Sylow subgroups of nilpotency class at most two. If the Sylow $2$-subgroups of $G$ are isomorphic to $\Q_8$, then $G$ is an $\IYB$-group.
\end{athm}


Recall that, for a skew left brace $(B,+,\cdot)$, the multiplicative group $(B,\cdot)$ acts on the additive group $(B,+)$ by automorphisms: for every $a\in B$, the map $\lambda_a \colon B \rightarrow B$, given by $\lambda_a(b) = -a + ab$, is an automorphism of $(B,+)$ and the map $\lambda\colon (B,\cdot) \rightarrow \Aut(B,+)$ which sends $a \mapsto \lambda_a$ is a group homomorphism (\cite{Rump2007}). Thus, for every $a,b\in B$ it holds
\begin{equation}
\label{eq:producte-lambda} a b = a+ \lambda_a(b).
\end{equation}
We denote $\Ker(B):= \ker \lambda $, the \emph{kernel of the group action}, and 
$$\Soc(B):=\Ker(B) \cap \Z(B,+)=\{a \in B \mid ab=a+b=b+a, \forall b \in B\}$$
\emph{the socle of $B$}, which is a subbrace of $B$.


\begin{definition}[{\cite[Definition~4]{BEJP2023}}]
\begin{enumerate}
\item Let $A$ be a group and let $(B,+,\cdot)$ be a skew left brace . We say that $A$ \emph{acts on the skew left brace} $(B,+,\cdot)$ if there is an action of $A$ on the set $B$ such that $a \cdot(g+h) = a\cdot g + a\cdot h$ and $a\cdot(gh) = (a\cdot g)(a\cdot h)$, for all $g, h \in B$, that is, if the action of $A$ on the set $B$ is actually an action of $A$ on the group $(B,+)$ and an action of $A$ on the group  $(B,\cdot)$.

\item An action of a group $A$ on an $\XX\YB$-group $(G, \cdot)$, for which it is understood that the associated skew left brace of $\XX$-type  is $(G, {+}, {\cdot})$, is said to be \emph{equivariant} if $A$ acts on the skew left brace $(G,+,\cdot)$. In this case, we say that $G$ has an \emph{$A$-equivariant} $\XX\YB$-structure.
\end{enumerate}
\end{definition}


The proof of Theorem~\ref{thm-Q-8}
heavily depends on two extensions of \cite[Theorem~A]{MengBallesterEstebanFuster2021} and \cite[Theorem~A]{BEJP2023}, which are  of independent interest. They hold for arbitrary (non-necessarily finite) groups.


\begin{theorem}\label{tA}
Let $\mathfrak{X}$ be a class of groups that is closed under taking quotients and direct products, and let the group $G = NH$ be the product of subgroups $N$ and $H$ with $N \trianglelefteq G$ and $N \cap H \leq \Z(N)$. Suppose that $N$ and $H$ are both $\mathfrak{X}\YB$-groups with, respectively, associated skew left braces $(N,+, \cdot)$, $(H,+, \cdot)$ satisfying the following conditions:

\begin{enumerate}
    \item $N \cap H \leq \operatorname{Ker}(N) \cap \operatorname{Ker}(H)$.
    \item $(N \cap H, +) \leq \Z(N, +) \cap \Z(H, +)$.
    \item The action of $H$ on $N$ by conjugation in $G$ is equivariant.
\end{enumerate}
Then, $G$ is an $\mathfrak{X}\YB$-group such that $\operatorname{Soc}(N) \operatorname{C}_{\operatorname{Soc}(H)}(N) \subseteq\operatorname{Soc}(G)$. Furthermore, the associated skew left braces $(N,+, \cdot)$ and $(H,+, \cdot)$ are subbraces of the associated skew left brace $(G,+, \cdot)$.
\end{theorem}

\begin{proof}
Most part of results follows from ~\cite[Theorem~A]{BEJP2023} and we only have to show that $\Soc(N)\C_{\Soc(H)}(N) \subseteq \Soc(G)$.
Note that, in the proof of~\cite[Theorem~A]{BEJP2023}, we get that $G=NH$ is an $\mathfrak{X}\YB$-group with the addition "+" defined by
$n_1h_1+n_2h_2=(n_1+n_2)(h_1+h_2),~~~n_1,n_2 \in N, h_1,h_2 \in H$
and  $\Soc(N)\C_{\Soc(H)}(N) \subseteq \Ker(N) \C_{\Ker(H)}(N) \subseteq \Ker(G)$.  By the definition of the addition, we observe that $\Z(N,+)\Z(H,+) \subseteq \Z(G,+)$, which implies that $\Soc(N)\C_{\Soc(H)}(N) \subseteq \Z(N,+)\Z(H,+) \subseteq \Z(G,+)$, as desired.
\end{proof}
%\begin{theorem}
%Let the group $G = NH$ be the product of  subgroups $N$ and $H$ with  $N\trianglelefteq G$ and $N\cap H \leq Z(N)$. Suppose that $N$ and $H$ are both IYB-groups with, respectively,  associated left braces $(N,+,\cdot)$, $(H,+,\cdot)$ satisfying the following conditions:

%\begin{enumerate}
%\item \label{cond1} $N \cap H \leq \Ker(N)\cap \Ker(H)$.
%\item \label{cond3} $(N,+,\cdot)$ is $H$-equivariant with respect the action of $H$ on $N$ by conjugation in $G$.
%\end{enumerate}
%Then $G$ has an IYB-structure $(G, +, \cdot)$ such that $\Ker(N)\C_{\Ker(H)}(N) \leq \Ker(G)$. Furthermore, $(N,+,\cdot)$ and $(H,+,\cdot)$ are subbraces of $(G,+,\cdot)$.

%If, moreover, $A$ is a group acting on $(G, \cdot)$ and the IYB-structures of $N$ and $H$ are $A$-equivariant, then $(G, +, \cdot)$  is an $A$-equivariant structure.
%\end{theorem}

%The proof of Theorem~\ref{thm-Q-8}
%heavily depends on the following extension (or consequence) of Theorem~\ref{tA} and , which is of independent interest. It is worth pointing out here that Theorem~\ref{thm-main} holds for arbitrary (non-necessarily finite) groups.

\begin{athm}\label{thm-main}\label{thm-skew-nilpotent-type}
Let $\mathfrak{X}$ be a class of groups that is closed under taking quotients and direct products, and let a group $A$ act on a group $(G, \cdot)$. Suppose that $G=A_1A_2\cdots A_k$ is a product of some $A$-invariant subgroups $A_i$,  $1 \leq i \leq k$. Write $H_0=G$, $H_i=A_{i+1}\cdots A_k$, for each $i \in \{1, \cdots,k-1\}$, and $A_{k+1} = H_k = 1$.
Suppose that $A_i$ has an $A$-equivariant $\XX\YB$-structure
$(A_i, +, \cdot)$
for each $i \in \{1, \cdots,k\}$, satisfying
the following conditions for each $i \in \{1, \cdots, k\}$:
\begin{itemize}
\item[\emph{(1)}] $A_i$ is normalised by $H_{i}$;
\item[\emph{(2)}] $(A_i, +, \cdot)$ is an $H_i$-equivariant $\XX\YB$-structure of $A_i$ with the respect to the action by conjugation of $H_i$ on $A_i$: ${}^ha=hah^{-1}$ for $h \in H_i, a \in A_i$;
\item[\emph{(3)}] $A_1A_2\cdots A_i \cap H_i \leq \C_G(A_iA_{i+1})$;
\item[\emph{(4)}] $\Z(A_i)\leq \Soc(A_i)$.
\end{itemize}
Then $G$ has an $A$-equivariant $\XX\YB$-structure $(G, +, \cdot)$ satisfying the following property:

(*) Suppose that $a=a_1a_2\cdots a_k\in \Z(A_1)\Z(A_2)\cdots \Z(A_k)$, where $a_i\in \Z(A_i)$, and $a_ia_{i+1}\cdots a_k\in \C_G(A_{i-1})$, for each $i=1, \cdots,k$. Then  $a \in \Soc(G)$.

Furthermore, $(A_i,+,\cdot)$ is a subbrace of $(G,+,\cdot)$ for every $i \in \{1, \cdots,k\}$.
\end{athm}

\begin{proof}
We argue by induction on $k$. It is clear that we may assume that  $k\geq 2$. It is rather clear that $H_1 = A_2\cdots A_{k}$ satisfies the hypotheses of the theorem. Then, by induction, $H_1$ has an $A$-equivariant $\XX\YB$-structure $(H_1,+,\cdot)$ satisfying Property~(*) and $(A_i,+,\cdot)$ are subbraces of $(H_1,+,\cdot)$ for every $i \in \{2, \cdots,k\}$.

Furthermore, $G = A_1H_1$, $H_1$ normalises $A_1$ and the $A$-equivariant $\XX\YB$-structure
$(A_1,+,\cdot )$ is $H_1$-equivariant with the respect to the action by conjugation of $H_1$ on
$A_1$ by Statement~(2). Let  $a_1=a_2\cdots a_{k} \in A_1\cap H_1$, where
$a_i \in A_i$,  $1\leq i \leq k$. Then $(a_1)^{-1}a_2\cdots a_k=1$. If $i \in \{2, \cdots, k-1\}$, then $(a_1)^{-1}a_2\cdots a_{i-1} a_i=(a_{i+1}\cdots a_{k})^{-1} \in A_1\cdots A_i\cap A_{i+1}\cdots A_{k}\leq \C_G(A_iA_{i+1})\leq \C_G(A_{i})$ and $(a_1)^{-1}a_2\cdots a_{i-1} = (a_{i}\cdots a_{k})^{-1}\in \C_G(A_{i-1}A_{i})\leq \C_G(A_{i})$ by Statement~(3). Therefore, $a_{i}\in \Z(A_{i})$. Now, $a_1\cdots a_{i-1}=(a_ia_{i+1}\cdots a_{k})^{-1}\in \C_G(A_{i-1})$ and $a_ia_{i+1}\cdots a_{k}\in \C_{G}(A_{i-1})$. Since $H_1$ satisfies Property~(*), $a_2\cdots a_k\in \Soc (H_1)$. Note that  $A_1\cap H_1\leq \C_G(A_1A_2)\cap A_1\leq \Z(A_1)\leq \Soc (A_1)$ by Statement~(4). Consequently, $A_1 \cap H_1 \leq \Soc (A_1) \cap \Soc (H_1)$. Applying Theorem~\ref{tA}, $G$ has an $A$-equivariant $\XX\YB$-structure $(G, +, \cdot)$ such that $\Soc(A_1)\C_{\Soc (H_1)}(A_1)\leq \Soc (G)$ and $(A_1, +, \cdot)$ and $(H_1,+,\cdot)$ are subbraces of $(G,+,\cdot)$. Hence $(A_i,+,\cdot)$ is a subbrace of $(G,+,\cdot)$, for every $i \in \{1, \cdots,k\}$.

Assume that $a=a_1a_2\cdots a_k\in \Z(A_1)\Z(A_2)\cdots \Z(A_k)$, where $a_i\in \Z(A_i)$, and $a_ia_{i+1}\cdots a_k\in \C_G(A_{i-1})$, for each $i=1, \cdots,k$. Then $a_2\cdots a_k\in \C_G(A_1)$ and $a_2\cdots a_k\in \Soc (H_1)$. Thus $a_1a_2\cdots a_k\in \Soc(A_1)\C_{\Soc (H_1)}(A_1)\leq \Soc (G)$. This shows that the skew left brace $(G, +, \cdot)$ has Property~(*).
\end{proof}

%Let a group $A$ act on a group $(G, \cdot)$ and suppose that $G=A_1A_2\cdots A_k$ is a product of some $A$-invariant subgroups $A_i$,  $1 \leq i \leq k$. Let $H_0=G$, $H_i=A_{i+1}\cdots A_k$, for each $i \in \{1, \cdots,k-1\}$, and $A_{k+1} = H_k = 1$.
%Suppose that $A_i$ has an $A$-equivariant $\IYB$-structure
%$(A_i, +, \cdot)$
%for each $i \in \{1, \cdots,k\}$, satisfying
%the following conditions for each $i \in \{1, \cdots, k\}$:
%\begin{itemize}
%\item[\emph{(1)}] $A_i$ is normalised by $H_{i}$;
%\item[\emph{(2)}] $(A_i, +, \cdot)$ is an $H_i$-equivariant IYB-structure of $A_i$ with the respect to the action by conjugation of $H_i$ on $A_i$: ${}^ha=hah^{-1}$ for $h \in H_i, a \in A_i$;
%\item[\emph{(3)}] $A_1A_2\cdots A_i \cap H_i \leq \C_G(A_iA_{i+1})$;
%\item[\emph{(4)}] $\Z(A_i)\leq \Ker(A_i)$.
%\end{itemize}
%Then $G$ has an $A$-equivariant IYB-structure $(G, +, \cdot)$ satisfying the following property:

%(*) Suppose that $a=a_1a_2\cdots a_k\in \Z(A_1)\Z(A_2)\cdots \Z(A_k)$, where $a_i\in \Z(A_i)$, and $a_ia_{i+1}\cdots a_k\in \C_G(A_{i-1})$, for each $i=1, \cdots,k$. Then  $a \in \Ker (G)$.

%Furthermore, $(A_i,+,\cdot)$ is a subbrace of $(G,+,\cdot)$ for every $i \in \{1, \cdots,k\}$.


%\begin{proof}

%We argue by induction on $k$. It is clear that we may assume that  $k\geq 2$. It is rather clear that $H_1 = A_2\cdots A_{k}$ satisfies the hypotheses of the theorem. Then, by induction, $H_1$ has an $A$-equivariant IYB-structure $(H_1,+,\cdot)$ satisfying Property~(*) and $(A_i,+,\cdot)$ are subbraces of $(H_1,+,\cdot)$ for every $i \in \{2, \cdots,k\}$.

%Furthermore, $G = A_1H_1$, $H_1$ normalises $A_1$ and the $A$-equivariant IYB-structure
%$(A_1,+,\cdot )$ is $H_1$-equivariant with the respect to the action by conjugation of $H_1$ on
%$A_1$ by Statement~(2). Let  $a_1=a_2\cdots a_{k} \in A_1\cap H_1$, where
%$a_i \in A_i$,  $1\leq i \leq k$. Then $(a_1)^{-1}a_2\cdots a_k=1$. If $i \in \{2, \cdots, k-1\}$, then $(a_1)^{-1}a_2\cdots a_{i-1} a_i=(a_{i+1}\cdots a_{k})^{-1} \in A_1\cdots A_i\cap A_{i+1}\cdots A_{k}\leq \C_G(A_iA_{i+1})\leq \C_G(A_{i})$ and $(a_1)^{-1}a_2\cdots a_{i-1} = (a_{i}\cdots a_{k})^{-1}\in \C_G(A_{i-1}A_{i})\leq \C_G(A_{i})$ by Statement~(3). Therefore, $a_{i}\in \Z(A_{i})$. Now, $a_1\cdots a_{i-1}=(a_ia_{i+1}\cdots a_{k})^{-1}\in \C_G(A_{i-1})$ and $a_ia_{i+1}\cdots a_{k}\in \C_{G}(A_{i-1})$. Since $H_1$ satisfies Property~(*), $a_2\cdots a_k\in \Ker (H_1)$. Note that  $A_1\cap H_1\leq \C_G(A_1A_2)\cap A_1\leq \Z(A_1)\leq \Ker (A_1)$ by Statement~(4). Consequently, $A_1 \cap H_1 \leq \Ker (A_1) \cap \Ker (H_1)$.

%Applying Theorem~\ref{tA}, $G$ has an $A$-equivariant IYB-structure $(G, +, \cdot)$ such that $\Ker(A_1)\C_{\Ker (H_1)}(A_1)\leq \Ker (G)$ and $(A_1, +, \cdot)$ and $(H_1,+,\cdot)$ are subbraces of $(G,+,\cdot)$. Hence $(A_i,+,\cdot)$ is a subbrace of $(G,+,\cdot)$, for every $i \in \{1, \cdots,k\}$.

%Assume that $a=a_1a_2\cdots a_k\in \Z(A_1)\Z(A_2)\cdots \Z(A_k)$, where $a_i\in \Z(A_i)$, and $a_ia_{i+1}\cdots a_k\in \C_G(A_{i-1})$, for each $i=1, \cdots,k$. Then $a_2\cdots a_k\in \C_G(A_1)$ and $a_2\cdots a_k\in \Ker (H_1)$. Thus $a_1a_2\cdots a_k\in \Ker(A_1)\C_{\Ker (H_1)}(A_1)\leq \Ker (G)$. This shows that the left brace $(G, +, \cdot)$ has Property~(*).
%\end{proof}

\emph{From now on, all groups considered will be finite.}


\section{Structural results: proof of Theorems~\ref{thm-Q-8} and~\ref{thm-Sylow-Q-8}}

This section contains the results needed for the proof of Theorem~\ref{thm-Q-8}. Our strategy is to show that a soluble group with Sylow subgroups of class at most two satisfies the structural hypotheses of Theorem~\ref{thm-main}. The elegant results developed by Hall to establish the foundations of the theory of soluble groups play a key role (\cite[Chapter~I]{DoerkHawkes1992}).

Let $G$ be a group. A \emph{Hall subgroup} of $G$ is a subgroup $H$ of $G$ such that
$(|G : H|, |H|) = 1$. If $p$ is a prime and $H$ is a Hall subgroup of $G$ such that $|G : H|$ is a power of $p$, then  $G = HP$ and $H \cap P = 1$ for every Sylow $p$-subgroup $P$ of $G$. For this reason, $H$ is sometimes called a \emph{Sylow $p$-complement} of $G$.

A celebrated result of Hall states that the existence of Sylow $p$-complements for each prime $p$ is a characteristic property of soluble groups, it is therefore not surprising that in the structural study of such groups Sylow complements plays a fundamental role (\cite[Theorems~I.3.3 and~I.3.5]{DoerkHawkes1992}).


A set $\mathcal{K}$ comprising the group $G$ together with exactly one Sylow $p$-complement of $G$ for each prime $p$ in the set $\pi(G)$ of all prime divisors of the order of $G$ is called a \emph{complemented basis} of $G$. The above result of Hall shows that $G$ is soluble if and only if $G$ has complemented basis and, in this case, $G$ has a transitive permutation representation when it acts by conjugation on the set of its complemented basis (\cite[Theorems~I.4.10 and~4.11]{DoerkHawkes1992}).


A \emph{system normaliser} of a soluble group $G$ is a stabiliser of this representation, that is, if  $\mathcal{K}$ be a complement basis of $G$, the system normaliser of $\mathcal{K}$ is  the intersection of the normalisers of every element of $\mathcal{K}$. It is denoted by $\N_G(\mathcal{K})$. It is worth noting that the system normalisers of $G$ are just the $\mathcal{N}$-normalisers of $G$ for the saturated formation $\mathcal{N}$ of all nilpotent groups (\cite[Section~3]{DoerkHawkes1992}).


According to \cite[Theorem~I.5.2]{DoerkHawkes1992}, the system normalisers of a soluble group $G$ form a characteristic conjugacy class of subgroups of $G$.

Let $G$ be a non-necessarily soluble group and let $F$ be a soluble subgroup of $G$. If $\mathcal{T}$ is a complement basis of $F$, then $\N_G(\mathcal{T}) = \{g \in G : H = H^g$ for all   $H \in \mathcal{T}\}$ is called \emph{system normaliser of $\mathcal{T}$  relative to $G$} (\cite{Hall1940}). Note that if $F$ is a normal subgroup of $G$, then $G = F\N_G(\mathcal{T})$ by the Frattini Argument.



%In this section, we will introduce some results about system normalisers of finite soluble groups. For more details, refer to \cite[Chapter IV]{Huppert2025}.

%Let $G$ be a finite group and $\pi(G)=\{p_1,\cdots, p_r\}$ denotes the set all primes dividing the order of $G$.
%A set $\mathcal{K} = \{G_{p_1'}, \ldots, G_{p_r'}\}$ is a called a complement system of $G$ if $G_{p_i'}$ is a complement of a Sylow $p_i$-subgroup of $G$.

%\begin{lemma}\label{lem-complement-system}
%Let $G$ be a group. Then
%\begin{itemize}
%\item[\emph{(1)}] $G$ is soluble if and only if $G$ has a complement system;
%\item[\emph{(2)}] Assume $G$ is soluble. Let $\mathcal{K} = \{G_{p_i'} \mid i=1,\ldots,r\}$
%and $\mathcal{K}^\ast = \{G_{p_i'}^\ast\mid i = 1,\ldots,r\}$ be complement systems of $G$. Then there exists an
%element $g \in G$ such that $G_{p_i'}^\ast = G_{p_i'}^g$ for all $i = 1,\ldots,r$.
%\end{itemize}
%\end{lemma}
%\begin{proof}
%See~\cite[Chapter IV, Theorem~2.3, Theorem~2.4]{Huppert2025}.
%\end{proof}

%Here, a more generalized notion of system normalisers due to P. Hall~\cite{Hall1940} is required.
%Let $F$ be a subgroup of $G$ and let $\mathcal{T}$ be a complement system of $F$. The subgroup
%$$\N_G(\mathcal{T}):=\bigcap_{X \in \mathcal{T}}\N_G(X)$$
%is the intersection of all normalisers of the subgroups $X$ in $\mathcal{T}$, which is called a system normaliser of $F$ relative to $G$.  In particular, $\N_F(\mathcal{T})$ (or simply denote by $\N(\mathcal{T})$) is call a system normaliser of $F$.

The following lemma is a slight generalization of \cite[Theorem~I.5.8]{DoerkHawkes1992} and we give a proof here for completeness.


\begin{lemma}\label{lem-system-normaliser}
Let $G$ be a group and let $F$ be a normal soluble subgroup of $G$.
Let $\mathcal{K}=\{F_{p'} \mid p \in \pi(F)\}$ be a complement basis of $F$. Let $N$ be a normal subgroup of $G$ contained in $F$. Then $\N_G(\mathcal{K})N/N$ is the system normaliser of the complement basis $\mathcal{K}N/N=\{F_{p'}N/N \mid p \in \pi(F/N)\}$ of $F/N$ relative to $G/N$.
\end{lemma}
\begin{proof}
Without loss of generality, we may assume that $N$ is a minimal normal subgroup of $G$. Then $N$ is an elementary abelian $q$-subgroup for some prime $q$.
As $\N_G(\mathcal{K})$ normalises each $F_{p'}$ in $\mathcal{K}$, $\N_G(\mathcal{K})N/N$ normalises each $F_{p'}N/N$ in $\mathcal{K}^\ast$. Hence
$\N_G(\mathcal{K})N/N \leq \N_{G/N}(\mathcal{K}^\ast).$ Conversely, for each $gN \in \N_{G/N}(\mathcal{K}^\ast)$ with $g \in G$,
$F_{p'}^gN=F_{p'}N$ for each prime $p$.
Assume that $p = q$. Then $F_{q'}^g$ and $F_{q'}$ are two $q'$-Hall subgroups of the soluble group $F_{q'}^g N=F_{q'}N$, there exists $n \in N$ such that $F_{q'}^{gn}=F_{q'}$. Write $x=gn$. Then $gN = xN$ normalises $F_{q'}N/N$. Assume that $p \neq q$. Then $N \leq F_{p'}^g \cap  F_{p'}$ and $F_{p'}^g=F_{p'}$. Then $gN$ normalises $F_{p'}N/N$. Therefore $gN \in \N_G(\mathcal{K})N/N$.
and $\N_{G/N}(\mathcal{K}^\ast) \leq \N_G(\mathcal{K})N/N$, as desired.
\end{proof}

%Recall that the \emph{nilpotent residual} of a group $G$, denoted by $G^{\mathcal{N}}$, the smallest normal subgroup of $G$ with nilpotent quotient.


%A group $G$ is said to be $\Q_8$-\emph{free} if $G$ has no section isomorphic to the quaternion group $Q_8$ of order $8$.


\begin{corollary}\label{nilres}
Let $G$ be a group and let $F$ be a normal soluble subgroup of $G$. If $D$ is a relative system normaliser of $F$ with respect to $G$, then $G = F^{\mathcal{N}}D$.
\end{corollary}

\begin{proof}
We have that $G = FD$ by the above observation. On the other hand, $H = D \cap F$ is a system normaliser of $F$ and, by Lemma~\ref{lem-system-normaliser}, $HF^{\mathcal{N}}/F^{\mathcal{N}}$ is a system normaliser of the nilpotent group $F/F^{\mathcal{N}}$. Since every $p$-complement of a nilpotent group is normal, it follows that $HF^{\mathcal{N}}/F^{\mathcal{N}} = F/F^{\mathcal{N}}$. Hence $F = F^{\mathcal{N}}H$ and so $G = F^{\mathcal{N}}D$.
\end{proof}


%The following result is a useful splitting theorem due to Carter and Hawkes \cite{CarterHawkes1967}.
%\begin{lemma}\label{lem-nilpotent-residual}
%Let $G$ be a soluble group such that $G^{\mathcal{N}}$ is Abelian. Then $G^{\mathcal{N}}$ is complemented in $G$ and all complements are conjugated in $G$. The complements are system normalisers of $G$.
%\end{lemma}
%\begin{proof}
%This follows from \cite[Theorem~5.15]{CarterHawkes1967} when applied to the formation consisting of all nilpotent groups.
%\end{proof}

%We present here a slight generalized version of Lemma~\ref{lem-nilpotent-residual} under the coprime action.
Suppose that a group $A$ acts on a group $G$ via automorphisms. Such action is called \emph{coprime} if $(|A|,|G|)=1$. According to \cite[8.2.6~(a)]{KurzweilStellmacher2004}, we have:
\begin{lemma}\label{lem-coprime-action}
Assume that the action of $A$ on a soluble group $G$ is coprime. For each prime $p$, there exists an $A$-invariant
$p$-complement of $G$.
\end{lemma}
%\begin{proof}
%It is a consequence of \cite[8.2.6~(a)]{KurzweilStellmacher2004}.
%\end{proof}

\begin{lemma}\label{lem-complement}
Assume that the action of $A$ on a soluble group $G$ is coprime. Let $F$ be a normal $A$-invariant subgroup of $G$. Then there exists a complement basis $\mathcal{K}$ of $F$ such that
\begin{itemize}
\item[(1)] $\mathbf{N}_G(\mathcal{K})$ is $A$-invariant;
\item[(2)] $G=F^{\mathcal{N}}\mathbf{N}_G(\mathcal{K})$;
\item[(3)] $F^{\mathcal{N}} \cap  \mathbf{N}_G(\mathcal{K}) \leq (F^{\mathcal{N}})'$.
\end{itemize}
\end{lemma}
\begin{proof}
Since $F$ is $A$-invariant and $F$ is soluble,  it follows from Lemma~\ref{lem-coprime-action} that there exists an $A$-invariant $p$-complement $F_{p'}$ of $F$ for each $p \in \pi(F)$. Set $\mathcal{K}=\{F_{p'} \mid  p \in \pi(F)\}$. Then $\mathcal{K}$ is a complemented basis of $F$. Let $H = \mathbf{N}_G(\mathcal{K})$ be the relative normaliser of $\mathcal{K}$ relative to $G$. Clearly $H$ is $A$-invariant as $F_{p'}$ is $A$-invariant for each $p \in \pi(F)$, and Statement~(1) holds.
Statement~(2) follows from Corollary~\ref{nilres}. Write $N=F^{\mathcal{N}}$ and $D = H \cap F$. Then $D$ is a system normaliser of $\mathcal{K}$ and $DN'/N'$ is the system normaliser of
$\mathcal{K}N'/N'$ of $F/N'$ by Lemma~\ref{lem-system-normaliser}. Applying Corollary~\ref{nilres}, $F/N' = (N/N')(DN'/N')$. Since $N/N'$ is the nilpotent residual of $F/N'$ and it is abelian, we can apply \cite[Theorems~IV.5.18 and~V.4.2]{DoerkHawkes1992}, to conclude that $N/N' \cap DN'/N' = 1$, that is, $N \cap D \leq N'$. Hence $N \cap H \leq N'$ and Statement~(3) holds.
\end{proof}


The following result is the key lemma to deal with $Q_8$-free groups. It is \cite[Theorem~4]{SeitzWright1969} for the formation $\mathcal{N}$ of all nilpotent groups.
\begin{lemma}\label{lem-Q8-free}
Let $G$ be a soluble group such that $G^{\mathcal{N}}$ is $\Q_8$-free. Then $G^{\mathcal{N}}$ has an abelian Sylow $2$-subgroup.
\end{lemma}



Lemma~\ref{lem-complement} leads to an interesting characterisation of solubility. As a consequence of the well-known theorem of Kegel and Wielandt if a group $G=A_1A_2\cdots A_k$ is the product of pairwise permutable nilpotent subgroups $A_1,A_2,\cdots,A_k$, then $G$ is soluble (\cite[Corollary~2.4.4]{AFG}). We can say even more.

\begin{definition}
An \emph{$\mathcal{N}$-decomposition} of a group $G$ is a set of nilpotent subgroups $A_1,A_2,\cdots,A_k$, $k\geq 2$, of $G$ satisfiying the following properties:
\begin{enumerate}
\item $G=A_1A_2\cdots A_k$;
\item $A_i$ is normalised by $A_j$ for each $1\leq i <j\leq k$.
\item $(A_1\cdots A_i) \cap (A_{i+1} \cdots A_k) \leq A_1' \cdots A_i'$ for each $1 \leq i \leq k-1$.
\end{enumerate}
\end{definition}

\begin{theorem}\label{Ndec}
Assume that the action of a group  $A$ on a group $G$ is coprime. Then $G$ is soluble if and only if $G$ has an $\mathcal{N}$-decomposition $\{A_1,A_2,\cdots,A_k\}$  composed of $A$-invariant subgroups $A_i$, $1 \leq i \leq k$. If $G^{\mathcal{N}}$ is $\Q_8$-free, then $A_i'$ of odd order, for every $1 \leq i \leq k-1$.
\end{theorem}
\begin{proof}
If $G$ has an $\mathcal{N}$-decomposition, then $G$ is soluble by the above observation.

Assume now that the action of $A$ on the soluble group $G$ is coprime. We prove that $G$ has an $\mathcal{N}$-decomposition $\{A_1,A_2,\cdots,A_k\}$ composed of $A$-invariant subgroups (with $A_i'$ of odd order for each $1\leq i \leq k-1$ if $G^{\mathcal{N}}$ is $\Q_8$-free).


If $G$ is nilpotent, then $\{1, G \}$ is an $\mathcal{N}$-decomposition of $G$ composed of $A$-invariant subgroups. Hence we may suppose that
$G$ is not nilpotent. Then, $\F(G) \neq G$. Let $\Fit_2(G)$ be the normal subgroup of $G$ such that $\Fit_2(G)/\Fit(G) = \Fit(G/\Fit(G))$. By \cite[Theorem~A.10.6]{DoerkHawkes1992}, $\Fit_2(G) \neq \Fit(G)$. In particular,  $A_1=\Fit_2(G)^{\mathcal{N}} \neq 1$ and $A_1 \leq \Fit(G)$. Since $A_1$ is characteristic in $G$, it follows that $A_1$ is normal in $G$ and $A$-invariant.

Now  assume that $G^{\mathcal{N}}$ is $\Q_8$-free. Since $A_1 \leq G^{\mathcal{N}}$, it follows that $A_1$ is $\Q_8$-free. Applying Lemma~\ref{lem-Q8-free}, $A_1$ has an abelian Sylow $2$-subgroup. As $A_1$ is nilpotent, $A_{1}'$ is the direct product of the commutator subgroups of its all Sylow subgroups. Hence  $A_{1}'$ is of odd order.


By Lemma~\ref{lem-complement}, there exists an $A$-invariant system normaliser $H_1$ of $\Fit_2(G)$ relative to $G$ such that $G=A_1H_1$ and $A_1 \cap H_1 \leq A_1'$. Then $H_1$ is an $A$-invariant soluble proper subgroup of $G$ by Lemma~\ref{lem-complement}. If $H_1$ is nilpotent, then $\{A_1, A_2 = H_1\}$ is the desired $\mathcal{N}$-decomposition of $G$. Assume that $H_1$ is not nilpotent.

Write $H_0=G$. Assume that $i >1$ and $G$ has an $A$-invariant nilpotent subgroups $A_j$, $1 \leq j \leq i-1$, and subgroups $H_j$,  $0 \leq j \leq i-1$ of $G$ such that

\begin{enumerate}
\item $H_j = A_{j+1} \cdots A_{i-1}$;
\item $A_j=\Fit_2(H_{j-1})^{\mathcal{N}} \leq \Fit(H_{j-1})$;
\item $H_j$ is an $A$-invariant system normaliser of $\Fit_2(H_{j-1})$ relative to $H_{j-1}$;
\item $H_{j-1}=A_jH_j$  and $A_j \cap H_j \leq A_j'$
\end{enumerate}

Assume that $H_{i-1}$ is not nilpotent. Then $A_i = \Fit_2(H_{i-1})^{\mathcal{N}} \neq 1$,  is an A-invariant nilpotent normal subgroup of $H_{i-1}$. Let $H_j$ be a system normaliser of
$\Fit_2(H_{i-1})$ relative to $H_{i-1}$. Then $H_j$ is $A$-invariant, $H_{i-1}=A_iH_i$ and $A_i \cap H_i \leq A_i'$ by Lemma~\ref{lem-complement}.

Since $G$ is finite, there exists $k \geq 2$ such that $A_k = H_{k-1}$ is nilpotent. Then $G=A_1A_2\cdots A_k$ and $A_i$ is normalised by $A_j$ for each $1\leq i <j\leq k$. We prove that $(A_1\cdots A_i) \cap (A_{i+1} \cdots A_k) \leq A_1' \cdots A_i'$ for each $1 \leq i \leq k-1$ by induction on $i$. It is clear that we may assume that $i >1$ and $B_{i-1} = (A_1\cdots A_{i-1}) \cap H_{i-1}= ( A_1' \cdots A_{i-1}')  \cap H_{i-1}$.

We have that  $B_{i-1} \unlhd H_{i-1}$. Denote with bars the images in the group $\overline{H_{i-1}}=H_{i-1}/B_{i-1}$ Write $F = \Fit_2(H_{i-1})$. Since $H_i$ is a system normaliser of $F$ relative to $H_{i-1}$, it follows that $\overline{H_{i}}$ is a system normaliser of $\overline{F}$ relative to $\overline{H_{i-1}}$ by  Lemma~\ref{lem-system-normaliser}. Hence $\overline{A_{i}} \cap \overline{H_{i}} \leq \overline{A_{i}}'$ by Lemma~\ref{lem-complement}. Thus $B_{i-1}A_i \cap H_i=B_{i-1}A_i' \cap H_i$

$$
\begin{aligned}
A_1A_2\cdots A_i \cap H_i&= A_1A_2\cdots A_i \cap H_{i-1}\cap H_i\\
&=(A_1A_2 \cdots A_{i-1} \cap H_{i-1})A_i \cap H_i\\
&=B_{i-1}A_i \cap H_i=B_{i-1}A_i' \cap H_i\\
&=(A_1'A_2'\cdots A_{i-1}' \cap H_{i-1})A_i' \cap H_i\\
&=A_1'A_2'\cdots A_{i-1}'A_i' \cap H_{i-1} \cap H_i\\
&=A_1'A_2'\cdots A_{i-1}'A_i' \cap H_i.
\end{aligned}
$$
Then $\{A_1,\cdots,A_k\}$ is an $\mathcal{N}$-decomposition of $G$.

Assume that $G^{\mathcal{N}}$ is $\Q_8$-free. Let $i \in \{1, \cdots, k-1\}$. Since $A_i \leq G^{\mathcal{N}}$, it follows that $A_i$ is $\Q_8$-free. Applying Lemma~\ref{lem-Q8-free}, $A_i$ has an abelian Sylow $2$-subgroup. Since $A_i$ is nilpotent, $A_{i}'$ is the direct product of the commutator subgroups of its Sylow subgroups. Hence $A_{i}'$ is of odd order.
\end{proof}


%Denote $A_2=\Fit_2(H_1)^{\mathcal{N}}$. By Lemma~\ref{lem-complement}, there is an $A$-invariant system normaliser $H_2$ of $\Fit_2(H_1)$ relative to $H_1$ such that $H_1=A_2H_2$ and $A_2 \cap H_2 \leq A_2'$.

%Write $H_0=G$. Arguing as above, we can construct by induction subgroups $A_i  (1 \leq i \leq k-1),H_i (0 \leq i \leq k-1)$ of $G$ such that
%\begin{itemize}
%\item[(a)] $A_i=\Fit_2(H_{i-1})^{\mathcal{N}} \leq \Fit(H_{i-1})$;
%\item[(b)] $H_i$ is an $A$-invariant system normaliser of $\Fit_2(H_{i-1})$ relative to $H_{i-1}$;
%\item[(c)] $H_{i-1}=A_iH_i$  and $A_i \cap H_i \leq A_i'$;
%\item[(d)] $H_{k-1}$ is nilpotent for some integer $k$;
%\end{itemize}
%We set $A_k=H_{k-1}$. It follows that $G=A_1A_2\cdots A_{k-1}H_{k-1}=A_1A_2\cdots A_{k-1}A_k$.
%and $H_i=A_{i+1}A_{i+2} \cdots A_k$ for $0 \leq i \leq k-1$, where $A_i$ is $A$-invariant and nilpotent. As $A_i=\Fit_2(H_{i-1})^{\mathcal{N}} \unlhd H_{i-1}$, $A_i$ is normalised by $A_{j} \leq H_i$ for $i<j$.



%Let $\XX$ be the set of all complement systems of $F$. Define the action of $G$ on $\XX$ via
%$$\mathcal{T}^g:=\{X^g \mid X \in \mathcal{T}\}~\text{for each $g \in G$ and each $\mathcal{T} \in \XX$.}$$
%We easily see that $G$ acts on $\XX$. By Lemma~\ref{lem-complement-system}(2), $F$ acts transitively on $\XX$. By Frattini's argument, $G=FH$, where $H=\mathbf{N}_G(\mathcal{K})$ is exactly the stabilizer of $\mathcal{K}$ in $G$. Note that $F \cap H$ is a system normaliser of $F$. It follows from \cite[Main Theorem~11.10]{Huppert2025} that $F \cap H$ covers all central chief factors of $F$. Hence $F \cap H$ covers $F/F^{\mathcal{N}}$, that is, $F \leq F^{\mathcal{N}}(F \capH)$.
%Then $G=F^{\mathcal{N}}H$.

%Write $N=F^{\mathcal{N}}$ and consider the quotient group $\overline{G}=G/N'$. We see that $\overline{G}=\overline{N} \cdot \overline{H}$. As $F=N(F \cap H)$, it implies that $\overline{F}=\overline{N} \cdot \overline{F \cap H}$.
%Note that $\overline{F}^{\mathcal{N}}=F^{\mathcal{N}}N'/N'=N/N'=\overline{N}$ is Abelian. As $F \cap H$ is a system normaliser of $F$,  by Lemma~\ref{lem-system-normaliser}, $\overline{F \cap H}$ is a system normaliser of $\overline{F}$. Applying Lemma~\ref{lem-nilpotent-residual}, $\overline{F \cap H}$ is a complement of $\overline{N}$ in $\overline{F}$. It implies that $N \cap (F \cap H)N' =N'$.  We deduce that
%$N'=N \cap (F \cap H)N'=(N \cap F \cap H)N'=(N \cap H)N'$, which implies that $N \cap H \leq N'$, as desired. The lemma is proved.
%\end{proof}










%\section{Applications}

%Recall that, for a group $G$, $\Fit(G)$ denotes the Fitting subgroup of $G$, which is the largest normal nilpotent subgroup of $G$. We also denotes by $\Fit_2(G)$ the preimage of $\Fit(G/\Fit(G))$ in $G$, so-called the second Fitting subgroup of $G$.

We shall make repeated use of the following two elementary lemmas.


\begin{lemma}\label{lem-class-two-property}
Let $G$ be a group with all Sylow subgroups of $G$ of nilpotency class at most two. Suppose that $N$ is a nilpotent subgroup of $G$. Then
\begin{itemize}
\item[\emph{(1)}] $N$ has nilpotency class at most two;
\item[\emph{(2)}] If $N$ is normalised by a subgroup $H$ of $G$, then $N' \cap H \leq\Z(\Fit(H))$.
\end{itemize}
\end{lemma}
\begin{proof}
As $N$ is nilpotent, $N=P_1 \times \cdots \times P_s$, where  $\pi(N)=\{p_1,\cdots,p_s\}$ and $P_i$ be a Sylow $p_i$-subgroup of $N$, $1 \leq i \leq s$. Since $P_i$ has nilpotent class at most $2$ for all $i$, we have that $N$ has nilpotency class at most two by \cite[Theorem~A.8.2]{DoerkHawkes1992} and Statement~$(1)$ holds.


Assume that $N$ is normalised by $H$. Then $N \cap H$ is a normal nilpotent subgroup of $H$ and so $N \cap H$ is contained in $\Fit(H)$, the Fitting subgroup of $H$. Write $N' \cap H=Q_1 \times \cdots \times Q_s$, where $Q_i$ is a Sylow $p_i$-subgroup of $N' \cap H$, $1 \leq i \leq s$.

Fix an index $i \in \{1, \cdots, s\}$. Since $N'=P_1' \times \cdots \times P_s'$, it follows that $Q_i \leq P_i'$ as $P_i'$ is the unique Sylow $p_i$-subgroup of $N'$.  Let $R_i$ be the Sylow $p_i$-subgroup of $\Fit(H)$. Then $Q_i \leq R_i$. Since $H$ normalises $N$, we have that $R_i$ normalises $P_i$. Hence $R_iP_i$ is a $p_i$-subgroup of $G$. Since $R_iP_i$ has nilpotency class at most two, it follows that $[Q_i,R_i]\leq [P_i', R_i]=1,$. Thus $Q_i \leq \Z(R_i) \leq \Z(\Fit(H))$. Consequently, $N' \cap H\leq \Z(\Fit(H))$ and Statement~(2) holds.
\end{proof}
%Now we prove Part~$(2)$.
%Since $N$ is normalised by $H$, $N \cap H \unlhd H$. The nilpotency of $N$ implies that $N' \cap H \subseteq N \cap H \subseteq \Fit(H)$. Write $\pi(N)=\{p_1,\cdots,p_s\}$ and let $P_i$ be a Sylow $p_i$-subgroup of $N$ for $1 \leq i \leq s$. As $N$ is nilpotent, it follows that $N=P_1 \times \cdots \times P_s$ and $N'=P_1' \times \cdots \times P_s'$. Since $Q_i \leq N' \cap H \leq N'$, then $Q_i \leq P_i'$ as $P_i'$ is the unique Sylow $p_i$-subgroup of $N'$,

%Set $N' \cap H=Q_1 \times \cdots \times Q_s$, where $Q_i$ is a Sylow $p_i$-subgroup of $N' \cap H$. As $Q_i \leq N' \cap H \leq N'$, we get that $Q_i \leq P_i'$ as $P_i'$ is the unique Sylow $p_i$-subgroup of $N'$. Hence $Q_i \leq P_i'$. Let $R_i$ be the Sylow $p_i$-subgroup of $\Fit(H)$ and clearly $Q_i \leq R_i$. Since $H$ normalises $N$, it implies that $R_i$ normalises $P_i$. Hence $R_iP_i$ is a $p_i$-subgroup, by hypothesis, $R_iP_i$ has nilpotency class at most two. Hence we see that
%$[Q_i,R_i]\leq [P_i', R_i]=1,$
%which implies that $Q_i \leq \Z(R_i) \leq \Z(\Fit(H))$. Thus $N' \cap H\leq \Z(\Fit(H))$, as desired.



\begin{lemma}\label{lem-Sylow}
Let a group $G=A_1A_2\cdots A_k$ be the product of some nilpotent subgroups $A_1,A_2,\cdots,A_k$ of $G$ such that $A_i$ is normalised by $A_j$ for $1 \leq i<j \leq k$. Let $p$ be a prime and let $P_i$ be the Sylow $p$-subgroup of $A_i$ for each $i$. Then $P_1P_2\cdots P_k$ is a Sylow $p$-subgroup of $G$.
\end{lemma}
\begin{proof}
We argue by induction on $k$. We can assume that $k > 1$  and the result is true for $k-1$. Note that $A=A_1A_2\cdots A_{k-1}$ is a subgroup of $G$ satisfying the hypotheses of the lemma. Hence $P=P_1P_2\cdots P_{k-1}$ is a Sylow $p$-subgroup of $A=A_1A_2\cdots A_{k-1}$. Since $A_i$ is nilpotent, $P_i$ is a characteristic subgroup of $A_i$. By hypothesis, $A_k$ normalises $A_i$ for $1 \leq i \leq k-1$. Hence $P_k \leq A_k$ normalises $P_i$ for $1 \leq i \leq k-1$. Consequently $P_k$ normalises $P=P_1P_2\cdots P_{k-1}$ and $PP_k$ is a $p$-subgroup of $G$. Let $|G|_p$be denote the largest power of $p$ dividing the order of $G$. Since $P \cap P_k$ is a $p$-subgroup of $A \cap A_k$, it follows that
$$|G|_p=\frac{|A|_p\cdot |A_k|_p}{|A \cap A_k|_p}~\text{divides}~\frac{|P|\cdot |P_k|}{|P \cap P_k|}=|PP_k|.$$
Hence $PP_k=P_1P_2\cdots P_k$ is a Sylow $p$-subgroup of $G$, as desired.
\end{proof}

The following lemmas turn out to be crucial to prove Theorem~\ref{thm-Q-8}.

\begin{lemma}\label{lem-nilpotency-class-2}
Assume that the action of a group  $A$ on a soluble group $G$ is coprime and all Sylow subgroups of $G$ have nilpotency class at most two. Let $\{A_1,A_2,\cdots,A_k\}$ be an $\mathcal{N}$-decomposition of $G$ composed of $A$-invariant subgroups. Then
$$A_1A_2\cdots A_i \cap A_{i+1}\cdots A_k \leq \C_G(A_iA_{i+1}) \text { for each } 1 \leq i \leq k-1.$$
%Then there exist nilpotent $A$-invariant subgroups $A_1,A_2,\cdots,A_k$ of $G$ such that
%\begin{itemize}
%\item[\emph{(1)}] $G=A_1A_2\cdots A_k$;
%\item[\emph{(2)}] $A_i$ is normalised by $A_j$ for $1\leq i <j\leq k$;
%\item[\emph{(3)}] $A_1A_2\cdots A_i \cap A_{i+1}\cdots A_k \leq \C_G(A_iA_{i+1})$ for each $1 \leq i \leq k-1$.
%\end{itemize}
%If $G^{\mathcal{N}}$ is $\Q_8$-free, then $A_i'$ is of odd order for each $1\leq i \leq k$.
\end{lemma}
\begin{proof}
%If $G$ is nilpotent, we can take $A_1=1,A_2=G$ and the result is trivial. Hence we may assume that
%$G$ is not nilpotent. Then, $\F(G) \neq G$. Let $\Fit_2(G)$ be the normal subgroup of $G$ such that $\Fit_2(G)/\Fit(G) = \Fit(G/\Fit(G))$. By \cite[Theorem~A.10.6]{DoerkHawkes1992}, $\Fit_2(G) \neq \Fit(G)$. In particular,  $A_1=\Fit_2(G)^{\mathcal{N}} \neq 1$ and $A_1 \leq \Fit(G)$. By Lemma~\ref{lem-complement}, there exists an $A$-invariant system normaliser $H_1$ of $\Fit_2(G)$ relative to $G$ such that $G=A_1H_1$ and $A_1 \cap H_1 \leq A_1'$. Denote $A_2=\Fit_2(H_1)^{\mathcal{N}}$. By Lemma~\ref{lem-complement}, there is an $A$-invariant system normaliser $H_2$ of $\Fit_2(H_1)$ relative to $H_1$ such that $H_1=A_2H_2$ and $A_2 \cap H_2 \leq A_2'$.

%Write $H_0=G$. Arguing as above, we can construct by induction subgroups $A_i  (1 \leq i \leq k-1),H_i (0 \leq i \leq k-1)$ of $G$ such that
%\begin{itemize}
%\item[(a)] $A_i=\Fit_2(H_{i-1})^{\mathcal{N}} \leq \Fit(H_{i-1})$;
%\item[(b)] $H_i$ is an $A$-invariant system normaliser of $\Fit_2(H_{i-1})$ relative to $H_{i-1}$;
%\item[(c)] $H_{i-1}=A_iH_i$  and $A_i \cap H_i \leq A_i'$;
%\item[(d)] $H_{k-1}$ is nilpotent for some integer $k$;
%\end{itemize}
%We set $A_k=H_{k-1}$. It follows that $G=A_1A_2\cdots A_{k-1}H_{k-1}=A_1A_2\cdots A_{k-1}A_k$.
%and $H_i=A_{i+1}A_{i+2} \cdots A_k$ for $0 \leq i \leq k-1$, where $A_i$ is $A$-invariant and nilpotent. As $A_i=\Fit_2(H_{i-1})^{\mathcal{N}} \unlhd H_{i-1}$, $A_i$ is normalised by $A_{j} \leq H_i$ for $i<j$. Hence we only have to prove that Statement~$(3)$ holds for the subgroups $A_1, A_2, \cdots, A_k$.

Write $H_i=A_{i+1}A_{i+2} \cdots A_k$ for $1 \leq i \leq k-1$. For each prime $p$, let $A_{i,p}$  be the Sylow $p$-subgroup of $A_i$, $1 \leq i \leq k$. By Lemma~\ref{lem-Sylow},
$G_p = A_{1,p}A_{2,p} \cdots A_{k,p}$ is a Sylow $p$-subgroup of $G$, and $A_{j,p}' \leq G_p'$ for each $j \in \{1, \cdots, k\}$.

Let $1 \leq i \leq k-1$ and write $B_i=A_1A_2\cdots A_i \cap H_i$. Let $B_{i,p}$ be a Sylow $p$-subgroup of $B_i$, Clearly $B_i \unlhd H_i$ as $H_i$ normalises $A_1A_2\cdots A_i$.

We prove the following statements by induction on $i$:
\begin{itemize}
\item[($\ast$-1)] $B_i=A_1'A_2' \cdots A_i' \cap H_i \leq \Z(\Fit(H_i)) \cap \C_G(A_{i})$;
\item[($\ast$-2)] $B_{i,p} \leq G_p'$ for each prime $p$.
\end{itemize}
If $i=1$, then
$B_1=A_1 \cap H_1 =A_1' \cap H_1$. Since $A_1$ is nilpotent and $A_{1,p}$ is the Sylow $p$-subgroup of $A_1$, we have that $A_{1,p}'$ is the Sylow $p$-subgroup of $A_1'$. Hence $B_{1,p} \leq A_{1,p}' \leq G_p'$ and ($\ast$-2) follows. Moreover, as $A_1$ is a nilpotent normal subgroup of $G$, it follows from Lemma~\ref{lem-class-two-property} that
$A_1' \cap H_1 \leq \Z(\Fit(H_1))$. Moreover, $A_1' \leq \Z(A_1)$ as $A_1$ is of nilpotency class $2$. Hence ($\ast$-1) and ($\ast$-2) holds for $i=1$.

Now we assume that $i > 1$ and ($\ast$-1) and ($\ast$-2) holds for $i-1$.
Then  $B_{i-1, p} \leq G_p'$ for each prime $p$ and
$B_{i-1}=A_1'A_2' \cdots A_{i-1}' \cap H_{i-1} \leq \Z(\Fit(H_{i-1})) \cap \C_G(A_{i-1}).$





%Write $F=\Fit_2(H_{i-1})$ and assume that $H_i=\mathbf{N}_G(\mathcal{K})$ for some complement system $\mathcal{K}$ of $F$.


%Note that $B_{i-1} \leq \Fit(H_{i-1}) \leq F$ and $B_{i-1} \unlhd H_{i-1}$. In the quotient group $\overline{H_{i-1}}=H_{i-1}/B_{i-1}$, we see that $\overline{F}=F/B_{i-1}$ and
%$\overline{F}^{\mathcal{N}}=\overline{F^{\mathcal{N}}}=\overline{A_{i}}=A_iB_{i-1}/B_{i-1}.$ Note that $H_i$ is a system normaliser of $F$ relative to $G$. By Lemma~\ref{lem-system-normaliser},  $\overline{H_{i}}$ is a system normaliser of $\overline{F}$ relative to $\overline{H_{i-1}}$. It follows from Lemma~\ref{lem-complement} that $\overline{A_{i}} \cap \overline{H_{i}} \leq \overline{A_{i}}',$ which implies that
%$B_{i-1}(B_{i-1}A_i \cap H_i)=B_{i-1}A_i \cap B_{i-1}H_i \leq B_{i-1}A_i'$. Hence $B_{i-1}A_i \cap H_i=B_{i-1}A_i' \cap H_i$. As $B_{i-1}=A_1'A_2' \cdots A_{i-1}' \cap H_{i-1}$ by induction, we get
%$$
%\begin{aligned}
%B_i=A_1A_2\cdots A_i \cap H_i&= A_1A_2\cdots A_i \cap H_{i-1}\cap H_i\\
%&=(A_1A_2 \cdots A_{i-1} \cap H_{i-1})A_i \cap H_i\\
%&=B_{i-1}A_i \cap H_i=B_{i-1}A_i' \cap H_i\\
%&=(A_1'A_2'\cdots A_{i-1}' \cap H_{i-1})A_i' \cap H_i\\
%&=A_1'A_2'\cdots A_{i-1}'A_i' \cap H_{i-1} \cap H_i\\
%&=A_1'A_2'\cdots A_{i-1}'A_i' \cap H_i.
%\end{aligned}
%$$
Since $B_{i-1}$ centralizes $A_i \leq \Fit(H_{i-1})$ and $A_i' \leq \Z(A_i)$ by hypothesis,
$B_{i-1}A_i'$ centralizes $A_i$ and so $B_i=B_{i-1}A_i' \cap H_i \leq \C_G(A_i)$.
Moreover, $B_i \leq B_{i-1}A_i'$ is nilpotent.
It follows from Lemma~\ref{lem-Sylow} that
$B_{i-1,p}A_{i,p}'$ is the unique Sylow $p$-subgroup of $B_{i-1}A_i'$ and so $B_{i,p} \leq B_{i-1,p}A_{i,p}'$. By induction, $B_{i-1,p} \leq G_p'$. Hence $B_{i,p} \leq B_{i-1,p}A_{i,p}' \leq G_p'$ and ($\ast$-2) holds for the case $i$.

Since $B_{i} \unlhd H_i$ and $B_i$ is nilpotent, we get $B_i \leq \Fit(H_i)$, moreover, $B_{i,p} \leq \Oo_p(H_i)$.
As $H_i=A_{i+1} \cdots A_{k}$, by Lemma~\ref{lem-Sylow}, $A_{i+1,p}A_{i+2,p} \cdots A_{k,p}$ is a Sylow $p$-subgroup of $H_i$. Clearly $\Oo_p(H_i) \leq A_{i+1,p}A_{i+2,p} \cdots A_{k,p} \leq G_p$.
Since $G_p$ is nilpotent of class two by hypothesis, $[B_{i,p}, \Oo_p(H_i)] \leq [G_p',G_p]=1$. Hence $B_{i,p} \leq \Z(\Oo_p(H_i))$, which implies that $B_i \leq \Z(\Fit(H_i))$ and ($\ast$-1) holds.

Since $A_{i+1} \leq \Fit(H_i)$, we get that
$A_1A_2\cdots A_i \cap A_{i+1}\cdots A_k \leq \C_G(A_iA_{i+1})$. The proof of the lemma is now complete.
\end{proof}

%Now we further assume that $G^{\mathcal{N}}$ is $\Q_8$-free.
%Clearly $\Fit_2(H_i)^{\mathcal{N}} \leq G^{\mathcal{N}}$ is $\Q_8$-free for each $i=0,1, \cdots,k-1$. Applying Lemma~\ref{lem-Q8-free} on $\Fit_2(H_i)$, we have that $A_{i+1}=\Fit_2(H_i)^{\mathcal{N}}$ has an Abelian Sylow $2$-subgroup. As $A_{i+1} \leq \Fit(H_i)$ is nilpotent, $A_{i+1}'$ is exactly the direct product of the commutator subgroups of its all Sylow subgroups. Hence  $A_{i+1}'$ is of odd order, as desired. The proof is complete.

%Using Theorem~\ref{thm-main} and Lemma~\ref{lem-nilpotency-class-2}, we now establish Theorem~\ref{thm-Q-8} and Corollary~\ref{cor-semidirect}.

The following known result is a consequence of Theorem~\ref{thm-main}.

\begin{lemma}[{\cite[Lemma~16]{CedoJespersRio2010}}]\label{nilp}
Every nilpotent group of nilpotency class at most two is an IYB-group.
\end{lemma}

\begin{proof}
Let $G$ be a nilpotent group of nilpotency class at most $2$. Then $G = A_1A_2\cdots A_k$ is the direct product of $\{ A_1, \cdots,A_k\}$  where $A_i$ is a Sylow $p_i$-subgroup of $G$, $1\leq i \leq k$, $\pi(G) = \{p_1,\cdots, p_k\}$. Assume that $k = 1$. Then $G$ is a $p$-group for some $p$, and $G$ is either cyclic or $G$ has an $\mathcal{N}$-decomposition of $G$  composed of abelian groups since $G/Z(G)$ is a direct product of cyclic subgroups. Since every abelian group is an IYB-group, it follows that $G$ is an IYB-group by Theorem~\ref{thm-main}.

Assume that $k \geq 2$. Then $\{ A_1, \cdots,A_k\}$ is an $\mathcal{N}$-decomposition of $G$ composed of IYB-groups. By Lemma~\ref{lem-nilpotency-class-2}, $G$ satisfies the hypotheses of Theorem~\ref{thm-main} for $A = 1$. Hence $G$ is an IYB-group.

\end{proof}

Applying~\cite[Corollary~4.2]{MengBallesterEstebanFuster2021}, we have:

\begin{lemma}\label{lem-odd-order-commutator}
Let $G$ be a nilpotent group of nilpotency class at most two with $|G'|$ odd. Then $G$ has an IYB-structure which is \(\Aut(G)\)-equivariant with respect to the natural action of \(\text{Aut}(G)\) on \( G \) and $\Z(G) \subseteq \Ker(G)$.
\end{lemma}



\begin{proof}[\textbf{\emph{Proof of Theorem~\ref{thm-Q-8}}}]
By Theorem~\ref{Ndec}, Lemmas~\ref{lem-nilpotency-class-2},~\ref{nilp} and~\ref{lem-odd-order-commutator}, $G$ has an $\mathcal{N}$-decomposition satisfying the statements of Theorem~\ref{thm-main} for $A = 1$. Hence $G$ is an IYB-group.
\end{proof}

In order to prove Theorem~\ref{thm-Sylow-Q-8}, we need the following IYB-structure on $\Q_8$.
\begin{lemma}\label{lem-IYB-Q-8}
Let $G=\Q_8$ and $A=\Aut(G)$. Then $G$ has an $\Oo^2(A)$-equivariant IYB-structure $(G,+,\cdot)$ with $\Z(G) \leq \Ker(G)$.
\end{lemma}
\begin{proof}
Let $G=\Q_8=\langle a,b \mid a^4=b^4=1, a^2=b^2, [a,b]=b^2\rangle$. Write $c=a^2=b^2 \in \Z(G)$.  Note that for every element $u \in G$, there exists the unique triple $(x,y,z) \in \mathbb{F}_2\oplus \mathbb{F}_2\oplus \mathbb{F}_2 $ such that $u=a^xb^yc^z$, where $\mathbb{F}_2=\{0,1\}$ is a field of order $2$. We define an addition on $G$ as
$$a^{x_1}b^{y_1}c^{z_1}+a^{x_2}b^{y_2}c^{z_2}= a^{x_1+x_2}b^{y_1+y_2}c^{z_1+z_2}.$$
Note that $1+1=0$ in the field $\mathbb{F}_2$. Hence $a^{x_1}a^{x_2}=a^{x_1+x_2}c^{x_1x_2}, b^{x_1}b^{x_2}=b^{x_1+x_2}c^{x_1x_2}$ and, more generally,
$$(a^{x_1}b^{y_1}c^{z_1})(a^{x_2}b^{y_2}c^{z_2})=a^{x_1+x_2} b^{y_1+y_2} c^{z_1+z_2+x_2y_1+y_1y_2+x_1x_2}$$
Then $(G,+,\cdot)$ is a left brace with $(G,+)$ elementary Abelian.
Note that $\lambda_c(u)=cu-c=a^xb^yc^{z+1}-c=a^xb^yc^z=u$ for each $u=a^xb^yc^z \in G$. Hence $\Z(G)=\langle c\rangle \leq \Ker(G)$.

Now we will show that this IYB-structure on $\Q_8$ is $\Oo^2(A)$-equivariant, that is,  $\Oo^2(A) \leq \Aut(G,+,\cdot)$. Since $A \cong S_4$, it suffices to prove that every element of $A$ with order $3$ belongs to $\Aut(G,+,\cdot)$. Note that every automorphism with order $3$ of $\Q_8$ permutates its three cyclic subgroups of order $4$. Hence we can easily list its all automorphism with order $3$:
$$
\begin{aligned}
&\alpha_1: a \mapsto b; b \mapsto ab;~~~~~~~~~~~\beta_1: a \mapsto ab; b \mapsto a;\\
&\alpha_2: a \mapsto b; b \mapsto (ab)c;~~~~~~\beta_2: a \mapsto (ab)c; b \mapsto a;\\
&\alpha_3: a \mapsto bc; b \mapsto ab;~~~~~~~~~\beta_3: a \mapsto (ab)c; b \mapsto ac;\\
&\alpha_4: a \mapsto bc; b \mapsto (ab)c;~~~~\beta_4: a \mapsto ab; b \mapsto ac;\\
\end{aligned}$$
Note that $\beta_i=\alpha_i^2$ for each $1 \leq i \leq 4$ and $\Oo^2(A)=\langle \alpha_i \mid 1 \leq i \leq 4 \rangle$.
Note that $(ab)^y=a^yb^y$ for $y=0,1$. We can check that $\alpha_i \in \Aut(G,+)$ for each $1 \leq i \leq 4$.

Since $\alpha_1(a^xb^yc^z)=b^x(ab)^yc^z=b^xa^yb^yc^z=a^yb^x[b^x,a^y]b^yc^z
=a^yb^xb^yc^{z+xy}=a^yb^{x+y}c^{xy}c^{z+xy}=a^yb^{x+y}c^z$ for $x,y,z \in \mathbb{F}_2$, it follows that
$$
\begin{aligned}
\alpha_1(a^{x_1}b^{y_1}c^{z_1}+a^{x_2}b^{y_2}c^{z_2})&=\alpha_1(a^{x_1+x_2}b^{y_1+y_2}c^{z_1+z_2})\\
&=a^{y_1+y_2}b^{x_1+x_2+y_1+y_2}c^{z_1+z_2}\\
&=a^{y_1}b^{x_1+y_1}c^{z_1}+a^{y_2}b^{x_2+y_2}c^{z_2}\\
&=\alpha_1(a^{x_1}b^{y_1}c^{z_1})+\alpha_1(a^{x_2}b^{y_2}c^{z_2}).
\end{aligned}
$$
Hence $\alpha_1 \in \Aut(G,+,\cdot)$. Similarly, for $x,y,z \in \mathbb{F}_2$, we have

$$
\begin{aligned}
\alpha_2(a^xb^yc^z)&=b^x(abc)^yc^z=a^yb^{x+y}c^{z+y}\\
\alpha_3(a^xb^yc^z)&=(bc)^x(ab)^yc^z=a^yb^{x+y}c^{z+x}\\
\alpha_4(a^xb^yc^z)&=(bc)^x(abc)^yc^z=a^yb^{x+y}c^{z+x+y}\\
\end{aligned}
$$
It can be directly verified that $\alpha_2,\alpha_3,\alpha_4 \in \Aut(G,+,\cdot)$ as well. Hence such IYB-structure on $\Q_8$ is $\Oo^2(A)$-equivariant and the proof is complete.
\end{proof}

Now we can prove Theorem~\ref{thm-Sylow-Q-8}.

\begin{proof}[\textbf{\emph{Proof of Theorem~\ref{thm-Sylow-Q-8}}}]
Since $G$ is soluble and all Sylow subgroups of $G$ have nilpotency class at most two, by Theorem~\ref{Ndec} and Lemma~\ref{lem-nilpotency-class-2}, $G$ has an $\mathcal{N}$-decomposition $\{A_1,A_2,\cdots,A_k\}$ such that
$A_1A_2\cdots A_i \cap A_{i+1}\cdots A_k \leq \C_G(A_iA_{i+1})$ for each $1\leq i \leq k-1$.
Write $H_i=A_{i+1}A_{i+2} \cdots A_k$ for each $1 \leq i \leq k-1$ and $H_k=1$.
We will apply Theorem~\ref{thm-main} for $A=1$. Hence we only have to show that $A_i$ has an $H_i$-equivariant IYB-structure $(A_i,+, \cdot)$ and $\Z(A_i) \leq \Ker(A_i)$ for each $i$.

Let $P_i$ be a Sylow $2$-subgroup of $A_i$ for each $i$ and write $P=P_1P_2\cdots P_k$, which is a
Sylow $2$-subgroup of $G$ by Lemma~\ref{lem-Sylow}.  Fixed some $i \in \{1,2,\cdots,k\}$. If $|P_i| \leq 2^2$, it implies that $P_i$ is Abelian. As $A_i$ is nilpotent, $A_i'$ is of odd order. It follows from Lemma~\ref{lem-odd-order-commutator} that $A_i$ has an $\Aut(A_i)$-equivariant IYB-structure $(A_i,+, \cdot)$ and $\Z(A_i) \leq \Ker(A_i)$, as desired.

Now assume that $|P_i| \geq 2^3$. It implies that $P_i=P$ as $|P|=2^3$ by hypothesis. Write $A_i=P \times W$, where $W$ is the Hall $2'$-subgroup of $A_i$. Considering the action of $H_i$ on $A_i=P \times W $, for each $j>i$, $P_j=P_j \cap P=P_j \cap P_i \leq A_i \cap H_j \leq \C_G(A_i)$. Hence $P_j \leq \Z(P)$, which implies that $P_{i+1}P_{i+2} \cdots P_k \leq \C_{H_i}(P)$. Note that $H_i/\C_{H_i}(P)$ is of order odd and acts faithfully on $P$. It follows from Lemma~\ref{lem-IYB-Q-8} that $P$ has an IYB-structure $(P,+,\cdot)$, which is $H_i/\C_{H_i}(P)$-equivariant, also $H_i$-equivariant, and $\Z(P) \leq \Ker(P)$.
Since $W$ is nilpotent of class two with odd order, by Lemma~\ref{lem-odd-order-commutator}, $W$ has an $H_i$-equivariant IYB-structure $(W,+,\cdot)$ with $\Z(W) \leq \Ker(W)$. It follows from Theorem~\ref{tA} that
$A_i$ has an $H_i$-equivariant IYB-structure $(A_i,+,\cdot)$ such that $\Z(A_i)=\Z(P)\Z(W) \leq \Ker(P)\Ker(W)=\Ker(A_i)$, as desired.
\end{proof}

\section{Some corollaries: skew braces of $\mathcal{N}$-type.}

In this section we show that \cite[Theorem~3.2,Theorem~3.5 ]{CedoOkninski2025} are natural consequences of our results. 


Let $(B,+,\cdot)$ a finite skew left brace of $\mathcal{N}$-type. Then $(B,+)$ is nilpotent and so it has a $\lambda$-invariant Hall $\pi$-subgroup for all set of primes $\pi$. In particular, $(B,\cdot)$ has Hall $\pi$-subgroups for all set of primes $\pi$. By \cite[Theorems~I.3.3 and~I.3.5]{DoerkHawkes1992}), $(B,\cdot)$ is soluble. Hence every finite $\mathcal{N}\YB$-group is soluble (see~\cite{Byott2015}).



\begin{corollary}
Let $G$ be a group with the Sylow tower property. Then $G$ is an $\mathcal{N}\YB$-group.
\end{corollary}
\begin{proof}
Since $G$ is a group with the Sylow tower property, we may assume that $G=A_1A_2\cdots A_k$ with $\pi(G)=\{p_1,p_2,\cdots,p_k\}$, where $A_i$ is a Sylow $p_i$-subgroup of $G$ such that $A_1A_2\cdots A_i \unlhd G$ and $A_1A_2\cdots A_i \cap A_{i+1}\cdots A_k=1$. Since $A_i$ is nilpotent, we can get $(A_i,+,\cdot)$ the trivial skew brace of nilpotent type, which implies that $A_i$ has fully equivariant $\mathcal{N}\YB$-structure $(A_i,+,\cdot)$ with $\Ker(A_i)=A_i$, moreover, $\Z(A_i)=\Z(A_i,+)=\Z(A_i,+) \cap \Ker(A_i)=\Soc(A_i)$.
Applying Theorem~\ref{thm-skew-nilpotent-type}, $G$ is an $\mathcal{N}\YB$-group.
\end{proof}

\begin{corollary}
Let $G$ be a soluble group with Sylow subgroups
of nilpotency class at most $2$. Then $G$ is a $\mathcal{N}\YB$-group.
\end{corollary}
\begin{proof}
We can apply Theorem~\ref{Ndec} and Lemma~\ref{lem-nilpotency-class-2} to conclude that 
$G$ has an $\mathcal{N}$-decomposition $\{A_1,\cdots,A_k\}$ such that
$A_1A_2\cdots A_i \cap A_{i+1}\cdots A_k \leq \C_G(A_iA_{i+1})$ for each $i$. Every nilpotent group $A_i$ has a trivial skew brace $(A_i,+,\cdot)$ with $\Z(A_i)=\Soc(A_i)$. By Theorem~\ref{thm-skew-nilpotent-type}, $G$ is an $\mathcal{N}\YB$-group.
\end{proof}



\textbf{Acknowledgment:} The first author is sponsored by Natural Science Foundation of Shanghai
(24ZR1422800) and National Natural Science Foundation of China (12471018). The second author is supported by the grant CIAICO/2023/007 from the Conselleria d'Educaci\'o, Universitats i Ocupaci\'o, Generalitat Valenciana and by the program "High-end Forcing Expert Program of China(H20240848)", Shanghai University.







\bibliographystyle{plain}
\bibliography{bibM}
\end{document}

------
Recall that a $1$-\emph{cocycle} of a $G$-module $V$ is a map $\pi: G \rightarrow V$ such that $\pi(gh)=\pi(g)+g\pi(h)$ for each $g,h \in G$. A group $G$ is an IYB-group if and only if there exists an $G$-module $V$ and a bijective $1$-cocycle $\pi: G \rightarrow V$. We call $(V,\pi)$ an IYB-structure of the IYB-group $G$ (see~\cite[Theorem~2.1]{CedoJespersRio2010}).

Let \( A \) be a group acting on an IYB-group \( G \) equipped with an IYB-structure \((V, \pi)\). For \( a \in A \) and \( g \in G \), we write \( {}^ag \) to denote the action of \( a \) on \( g \). Following~\cite{MengBallesterEstebanFuster2021}, the IYB-structure \((V, \pi)\) is called \textit{A-equivariant} if there exists a group action of \( A \) on \( V \) (denoted by \( a \cdot v \) for \( a \in A \) and \( v \in V \)) which satisfies the follwing compatibility condition:
\[
\pi({}^ag) = a \cdot \pi(g) \quad \text{for all } a \in A \text{ and } g \in G.
\]
Due to the bijectivity of \(\pi\), this action of \( A \) on \( V \) is uniquely determined by the action of \( A \) on \( G \) through the formula:
\[
a \cdot v = \pi \left( {}^a\pi^{-1}(v) \right), \quad \forall a \in A, v \in V.
\]

Furthermore, an IYB-structure \((V, \pi)\) is \( A \)-equivariant if and only if it is \( A/K \)-equivariant, where \( K = \text{Ker}(A~\text{on}~G) \) is the kernel of the action of \( A \) on \( G \). Note that the kernel of the action of $A$ on $G$ coincides the kernel of the action of $A$ on $V$, that is, $\text{Ker}(A~\text{on}~G)=\text{Ker}(A~\text{on}~V)$.
We say that an IYB-structure \((V, \pi)\) on \( G \) is \textit{fully equivariant} if it is \(\text{Aut}(G)\)-equivariant with respect to the natural action of \(\text{Aut}(G)\) on \( G \). In this case, \((V, \pi)\) automatically becomes \( A \)-equivariant for any group action of \( A \) on \( G \).
-----






-----
\section{IYB-structures: the proof of Theorem~\ref{thm-main}}
-----



In this section, we collect some results about IYB-structures of  IYB-groups which are needed to prove Theorem~\ref{thm-main}. Most of them appear in~\cite{MengBallesterEstebanFuster2021}, but we include them for the sake of completeness.


Our first lemma follows from \cite[Lemma~2.6]{MengBallesterEstebanFuster2021} and a direct calculation.

\begin{lemma}\label{lem-1-cocycle}
Let $G$ be a group and let \(\pi\) be a \(1\)-cocycle of a \(G\)-module \(V\). Suppose that \(g \in G\) and \(x \in \operatorname{Ker}(G \text{ on } V)\). Then
\begin{itemize}
\item[(1)] \(\pi(xg) = \pi(x) + \pi(g);\)
\item[(2)] $\pi(1)=0, \pi(x^{-1})=-\pi(x)$;
\item[(3)] \(\pi(gxg^{-1}) = g\pi(x).\)
\end{itemize}
\end{lemma}


The second lemma is a special case of \cite[Theorem~A]{MengBallesterEstebanFuster2021} for $A=1$ and $N \cap H=1$.

\begin{lemma}\label{lem-equivariant-IYB}
Let a group $G=NH$ be the product of subgroups $H$ and $N$ such that $N \unlhd G$ and $H \cap N=1$. Suppose that $H$ is an IYB-group and $N$ has an $A$-equivariant IYB-structure. Then $G$ is an IYB-group.
\end{lemma}


Applying~\cite[Corollary~4.2]{MengBallesterEstebanFuster2021}, we have:

\begin{lemma}\label{lem-odd-order-commutator}
Let $G$ be a nilpotent group of class at most two with $|G'|$ odd. Then $G$ has a fully equivariant IYB-structure $(V,\pi)$ with $\Z(G) \subseteq \Ker(G~on~V)$.
\end{lemma}

The next lemma is a consequence of \cite[Proposition~2.5]{MengBallesterEstebanFuster2021}.
\begin{lemma}\label{lem-nilpotency-IYB}
Let $G$ be a nilpotent group of class at most two. Then $G$ has a $\Aut_c(G)$-equivariant IYB-structure $(V,\pi)$ with $\Z(G) \subseteq \Ker(G~on~V)$.
\end{lemma}


The factorised group described in our next lemma has very good structural properties to generate IYB-structures.

\begin{lemma}\label{lem-elements}
Let $G=A_1A_2\cdots A_k$ be the product of the subgroups $A_i$,  $1 \leq i \leq k$, such that $A_i$ is normalised by $A_j$ for $1\leq i<j \leq k$. Suppose that
$$A_1A_2 \cdots A_i \cap A_{i+1}\cdots A_k \leq \C_G(A_iA_{i+1}),~i=1,2,\cdots,k-1.$$
Let $a_1a_2\cdots a_k=b_1b_2\cdots b_k$, where $a_i,b_i \in A_i$ for each $i$. Then
\begin{itemize}
\item[\emph{(1)}] For each $1\leq i \leq k-1$, we have that
$$\emph{\text{(1-1)}}~(b_1^{-1}a_1)(b_2^{-1}a_2)\cdots (b_i^{-1}a_i)=(b_{i+1}\cdots b_k)(a_{i+1}\cdots a_k)^{-1} \in \C_G(A_iA_{i+1});$$
$$\emph{\text{(1-2)}}~(a_{i+1}b_{i+1}^{-1})(a_{i+2}b_{i+2}^{-1})\cdots (a_kb_k^{-1})=(a_{1}\cdots a_{i})^{-1} (b_{1}\cdots b_{i}) \in \C_G(A_iA_{i+1});$$
\item[\emph{(2)}] $a_ib_i^{-1}, b_i^{-1}a_i \in \Z(A_i)$ for each $1 \leq i \leq k$; In particular if $b_1=b_2=\cdots=b_k=1$, then $a_i \in \Z(A_i).$
\item[\emph{(3)}] Let $h_i=a_{i+1}a_{i+2} \cdots a_{k}~(1\leq i \leq k-1)$ and $h_k=1$. Suppose that
$x_1x_2 \cdots x_k=1$ for some $x_i \in A_i, 1\leq i \leq k$. Then
$$({}^{h_1}x_1)\cdot ({}^{h_2}x_2) \cdots ({}^{h_k}x_k)=1. $$
\item[\emph{(4)}] Write $$\XX_j=\{(a_1,a_2 \cdots, a_j) \mid a_1a_2\cdots a_j=1, a_i \in A_i, \forall 1 \leq i \leq j\}$$
    for each $j=1,2,\cdots,k$. Then
    $$|A_1A_2\cdots A_j| \cdot |\XX_j|=|A_1|\cdot |A_2| \cdots |A_j|$$
for each $j=1,2,\cdots,k.$
\end{itemize}
\end{lemma}
\begin{proof}
(1) We will first prove $\text{(1-1)}$ by induction on $i$. For $i=1$,
$$b_1^{-1}a_1=(b_{2}\cdots b_k)(a_{2}\cdots a_k)^{-1} \in A_1 \cap A_{2}A_3\cdots A_k \leq \C_G(A_1A_{2}).$$
Now we assume that $\text{(1-1)}$ holds for $i-1$, that is,  $$(b_1^{-1}a_1)(b_2^{-1}a_2)\cdots (b_{i-1}^{-1}a_{i-1})=(b_{i}\cdots b_k)(a_{i}\cdots a_k)^{-1} \in \C_G(A_{i-1}A_{i}),$$
which implies that
$$b_i^{-1}(b_1^{-1}a_1)(b_2^{-1}a_2)\cdots (b_{i-1}^{-1}a_{i-1})a_{i+1}=(b_{i+1}\cdots b_k)(a_{i+1}\cdots a_k)^{-1}. $$
Note that $(b_1^{-1}a_1)(b_2^{-1}a_2)\cdots (b_{i-1}^{-1}a_{i-1}) \in \C_G(A_i)$ by induction, hence it centralizes $b_i^{-1} \in A_i$. As $A_{i+1}A_{i+2} \cdots A_k \leq G$,  we get that
$$(b_1^{-1}a_1)(b_2^{-1}a_2)\cdots (b_i^{-1}a_i)=(b_{i+1}\cdots b_k)(a_{i+1}\cdots a_k)^{-1}$$
belongs to $A_1A_2\cdots A_i \cap A_{i+1}\cdots A_k \leq \C_G(A_iA_{i+1})$, as desired.
We next show (1-2) by backward induction on $i$. For $i=k-1$, by hypothesis,
$$a_kb_k^{-1}=(a_{1}\cdots a_{k-1})^{-1} (b_{1}\cdots b_{k-1}) \in A_1A_2\cdots A_{k-1} \cap A_k \leq \C_G(A_{k-1}A_k).$$
We assume that it is true for $i+1$, that is,
$$(a_{i+2}b_{i+2}^{-1})(a_{i+3}b_{i+3}^{-1})\cdots (a_kb_k^{-1})=(a_{1}\cdots a_{i+1})^{-1} (b_{1}\cdots b_{i+1}) \in \C_G(A_{i+1}A_{i+2}),$$
which implies that
$$a_{i+1}(a_{i+2}b_{i+2}^{-1})(a_{i+3}b_{i+3}^{-1})\cdots (a_kb_k^{-1})b_{i+1}^{-1}=(a_{1}\cdots a_{i})^{-1} (b_{1}\cdots b_{i}).$$
Since $(a_{i+2}b_{i+2}^{-1})(a_{i+3}b_{i+3}^{-1})\cdots (a_kb_k^{-1}) \in \C_G(A_{i+1})$ by induction, which commutes with $b_{i+1}^{-1} \in A_{i+1}$, it follows that
$$(a_{i+1}b_{i+1}^{-1})(a_{i+2}b_{i+2}^{-1})\cdots (a_kb_k^{-1})=(a_{1}\cdots a_{i})^{-1} (b_{1}\cdots b_{i}),$$
which is in $A_1 A_2\cdots A_i \cap A_{i+1} \cdots A_k \leq
\C_G(A_iA_{i+1})$. Therefore Statement~$(1)$ is proved.
\\
\\
(2) Since that Equation (1-1) holds for $i=1,k-1$, we have that
$$b_1^{-1}a_1=(b_{2}\cdots b_k)(a_{2}\cdots a_k)^{-1} \in \C_G(A_1A_{2}) \cap A_1 \leq \Z(A_1);~\text{and}$$
$$(b_1^{-1}a_1)(b_2^{-1}a_2)\cdots (b_{k-1}^{-1}a_{k-1})=b_ka_k^{-1} \in \C_G(A_{k-1}A_{k}) \cap A_k \leq \Z(A_k).$$
Hence $b_i^{-1}a_i \in \Z(A_i)$ for $i=1$ and $k$. Moreover, for each $2\leq i \leq k-1$, it follows from Statement~(1) that
$$(b_1^{-1}a_1)(b_2^{-1}a_2)\cdots(b_{i-1}^{-1}a_{i-1}) (b_i^{-1}a_i) \in \C_G(A_iA_{i+1}) \leq \C_G(A_i);$$
and
$$(b_1^{-1}a_1)(b_2^{-1}a_2)\cdots(b_{i-1}^{-1}a_{i-1}) \in \C_G(A_{i-1}A_i) \leq \C_G(A_i).$$
Thus $b_i^{-1}a_i \in \C_G(A_i) \cap A_i=\Z(A_i)$ for $2\leq i \leq k-1$.  Hence we deduce that $b_i^{-1}a_i \in \Z(A_i)$ for each $1 \leq i \leq k$, and so $b_i^{-1}a_i=b^{-1}_i(a_ib_i^{-1})b_i \in \Z(A_i)$, as desired.
\\
\\
(3) In fact, we will prove
$$({}^{h_i}x_i)\cdot ({}^{h_{i+1}}x_{i+1}) \cdots ({}^{h_k}x_k)={}^{h_i}(x_ix_{i+1}\cdots x_k)$$
by backward induction on $i$. It is trivial for $i=k$ and we assume that
$$({}^{h_{i+1}}x_{i+1}) \cdots ({}^{h_k}x_k)={}^{h_{i+1}}(x_{i+1}\cdots x_k)$$
holds, where $i \leq k-1$.
Note that $h_i=a_{i+1}h_{i+1}$ and $$h_i^{-1}h_{i+1}=h_{i+1}^{-1}a_{i+1}^{-1}h_{i+1} \in A_{i+1}$$
since $h_{i+1}=a_{i+2}a_{i+3} \cdots a_{k} \in A_{i+2}A_{i+3}\cdots A_k$ normalises $A_{i+1}$ by hypothesis.
As $x_1x_2\cdots x_k=1$ and $A_1A_2\cdots A_i \leq G$ by hypothesis, we have
$x_{i+1}x_{i+2}\cdots x_k=(x_1x_2\cdots x_i)^{-1} \in A_1A_2\cdots A_i \cap A_{i+1} \cdots A_k \leq \C_G(A_iA_{i+1}).$
Hence $x_{i+1}x_{i+2}\cdots x_k$ commutes with $h_i^{-1}h_{i+1}$. Hence
$$
\begin{aligned}
({}^{h_{i}}x_{i})({}^{h_{i+1}}x_{i+1}) \cdots ({}^{h_k}x_k)&=({}^{h_{i}}x_{i}) \cdot {}^{h_{i+1}}(x_{i+1}x_{i+2}\cdots x_k)\\
&=h_ix_i(h_i^{-1}h_{i+1})(x_{i+1}x_{i+2}\cdots x_k)h_{i+1}^{-1}\\
&=h_ix_i(x_{i+1}x_{i+2}\cdots x_k)(h_i^{-1}h_{i+1})h_{i+1}^{-1}\\
&={}^{h_i}(x_ix_{i+1}x_{i+2}\cdots x_k),
\end{aligned}
$$
as desired. Consequently, $({}^{h_1}x_1)\cdot ({}^{h_{2}}x_{2}) \cdots ({}^{h_k}x_k)={}^{h_1}(x_1x_{2}\cdots x_k)=1$ and Statement~(3) holds.
\\
\\
(4) We prove Statement~(4) by induction on $j$. Clearly, we may assume that $j \geq 2$. Note that $A_1A_2 \cdots A_{j-1}$ is a subgroup of $G$. By induction we have
$$
\begin{aligned}
|A_1|\cdot |A_2| \cdots |A_{j-1}| \cdot |A_j|&=|A_1A_2\cdots A_{j-1}| \cdot |\XX_{j-1}| \cdot |A_j|\\
&= |A_1A_2\cdots A_{j-1}|\cdot |A_j| \cdot |\XX_{j-1}| \\
&=|A_1A_2\cdots A_j| \cdot |A_1A_2\cdots A_{j-1} \cap A_j| \cdot |\XX_{j-1}|
\end{aligned}
$$
Write $B=A_1A_2\cdots A_{j-1} \cap A_j$. We only have to show that
$|\XX_j|=|B| \cdot |\XX_{j-1}|$.
For $b \in A_j$, we consider the set
$$\XX_j(b)=\{(a_1,a_2,\cdots,a_{j-1},b) \mid a_1a_2\cdots a_{j-1}b=1, a_i \in A_i, \forall 1 \leq i \leq j-1\},$$
which is a subset of $\XX_j$.
 It is easy to see that $\XX_j(b) \neq \varnothing$ if and only if $b \in B$. Note that $\XX_j$ is the disjoint union of all $\XX_j(b)$ for $b \in B$. It implies that $|\XX_j|=\sum_{b \in B} |\XX_j(b)|.$

We claim that $|\XX_j(b)|=|\XX_{j-1}|$ for each $b \in B$. Fix an element $b \in B$, as $B=A_1A_2\cdots A_{j-1} \cap A_j$, we may assume that $b^{-1}=b_1b_2\cdots b_{j-1}$ for some $b_i \in A_i, 1\leq i \leq j-1$. Let
$\phi: \XX_j(b) \rightarrow \XX_{j-1}$ be a correspondence defined as follows:
$$\phi((a_1,a_2,\cdots, a_{j-1}, b))=(a_1b_1^{-1}, a_2b_2^{-1}, \cdots, a_{j-1}b_{j-1}^{-1}).$$
Since $a_1a_2\cdots a_{j-1}=b^{-1}=b_1b_2\cdots b_{j-1}$, by Statement~(1), it follows that $(a_1b_1^{-1})(a_2b_2^{-1})\cdots(a_{j-1}b_{j-1}^{-1})=1$. Hence

$$\phi((a_1,a_2,\cdots, a_{j-1}, b))=(a_1b_1^{-1}, a_2b_2^{-1}, \cdots, a_{j-1}b_{j-1}^{-1}) \in \XX_{j-1}.$$
Hence $\phi$ is actually a map. For each two elements $(a_1,a_2,\cdots,a_{j-1},b)$ and $(a_1',a_2',\cdots,a_{j-1}',b)$ in $\XX_j(b)$ with
$$\phi((a_1,a_2,\cdots,a_{j-1},b))=\phi((a_1',a_2',\cdots,a_{j-1}',b)),$$
we have that $a_ib_i^{-1}=a_i'b_i^{-1}$ for $1 \leq i \leq j-1$, and so $a_i=a_i'$. Hence $\phi$ is injective.  For each $(a_1,a_2,\cdots, a_{j-1}) \in \XX_{j-1}$,
we have $a_1a_2\cdots a_{j-1}=1$.
Write $l_0=1$, $l_i=a_1 a_2\cdots a_{i}, 1\leq i \leq j-1$ and note that $l_{j-1}=a_1 a_2\cdots a_{j-1}=1$. Then
$$l_i=a_1a_2\cdots a_{i}=(a_{i+1}\cdots a_{j-1})^{-1} \in A_1A_2\cdots A_i \cap A_{i+1} \cdots A_k$$
centralizes $b_i \in A_i$ by hypothesis.
As $a_{i}=l_{i-1}^{-1}l_{i}$ for $1\leq i \leq j-1$, we have that
$$
\begin{aligned}
&(a_1b_1)(a_2b_2)\cdots (a_{j-2}b_{j-2})(a_{j-1}b_{j-1})b\\
=&(l_1b_1)(l_1^{-1}l_2b_2)(l_2^{-1}l_3b_3) \cdots (l_{j-3}^{-1}l_{j-2}b_{j-2})(l_{j-2}^{-1}l_{j-1}b_{j-1})b\\
=&({}^{l_1}b_1)({}^{l_2}b_2)({}^{l_3}b_3) \cdots ({}^{l_{j-2}}b_{j-2}) b_{j-1}b\\
=&(b_1b_2\cdots b_{j-2}b_{j-1})b=b^{-1}b=1,
\end{aligned}
$$
which implies that $(a_1b_1,a_2b_2,\cdots,a_{j-1}b_{j-1}, b) \in \XX_j(b)$. Then
$$\phi((a_1b_1,a_2b_2,\cdots, a_{j-1}b_{j-1}, b))=(a_1, a_2, \cdots, a_{j-1}).$$
Hence $\phi$ is bijective and $|\XX_j(b)|=|\XX_{j-1}|$ for each $b \in B$. Consequently,
$$|\XX_j|=\sum_{b \in B} |\XX_j(b)|=|B||\XX_{j-1}|,$$
as desired. The proof of the lemma is now complete.
\end{proof}
\begin{proof}[\emph{\textbf{Proof of Theorem~\ref{thm-main}}}]
 Write
 $$X=\{(\pi_1(a_1), \pi_2(a_2), \cdots, \pi_k(a_k)) \mid a_1a_2 \cdots a_k=1, a_i \in A_i, \forall 1\leq i \leq k\},$$
 which is a subset of the $A$-module $U_1 \oplus U_2 \cdots \oplus U_k$.
\\
\\
\textbf{Claim 1}. $X$ is an $A$-submodule of $U_1\oplus U_2 \cdots \oplus U_k$.

Let $a_i,b_i \in A_i$ such that $a_1a_2\cdots a_k=1=b_1b_2\cdots b_k$. By Lemma~\ref{lem-elements}(2) and Hypothesis (4),  we have $a_i, b_i \in \Z(A_i) \subseteq \Ker(A_i~on~U_i)$. By Lemma~\ref{lem-1-cocycle}, $\pi_i(b_i^{-1}a_i)=\pi_i(b_i^{-1})+\pi_i(a_i)=\pi_i(a_i)-\pi_i(b_i)$. Hence
$$\begin{aligned}
&(\pi_1(a_1), \pi_2(a_2), \cdots, \pi_k(a_k))-(\pi_1(b_1), \pi_2(b_2), \cdots, \pi_k(b_k))\\
=& (\pi_1(a_1)-\pi_1(b_1), \pi_2(a_2)-\pi_2(b_2), \cdots, \pi_k(a_k)-\pi_k(b_k))\\
=& (\pi_1(b_1^{-1}a_1), \pi_2(b_2^{-1}a_2), \cdots, \pi_k(b_k^{-1}a_k)) \in X,
\end{aligned}
$$
as $(b_1^{-1}a_1)(b_2^{-1}a_2)\cdots (b_k^{-1}a_k)=1$ by Lemma~\ref{lem-elements}(1); moreover, for $a \in A$, we see that
$$
\begin{aligned}
&a(\pi_1(a_1), \pi_2(a_2), \cdots, \pi_k(a_k))\\
=&(a\pi_1(a_1), a\pi_2(a_2), \cdots, a\pi_k(a_k))\\
=& (\pi_1({}^aa_1), \pi_2({}^aa_2), \cdots, \pi_k({}^aa_k)) \in X,
\end{aligned}
$$
as $({}^aa_1)({}^aa_2)\cdots ({}^aa_k)={}^a(a_1a_2\cdots a_k)=1$. Hence $X$ is an $A$-submodule of $U_1\oplus U_2 \cdots \oplus U_k$, as claimed.
\\
\\
Consider now the corresponding quotient $A$-module $W=(U_1\oplus U_2 \cdots \oplus U_k)/X$. By Hypothesis~(2), $(U_i, \pi_i)$ is an $H_i$-equivariant IYB-structure on $A_i$, which determines an action of $H_i$ on $U_i$ such that
$hu=\pi_i(h\pi_i^{-1}(u)h^{-1})$ for $h \in H$ and $u \in U_i$.
For each $c=a_1a_2\cdots a_k \in G$, where $a_i \in A_i$, we write $h_i=a_{i+1}\cdots a_k$ for $1\leq i \leq k-1$ and $h_k=1$. We define a correspondence from $G \times W$ to $W$ given by
$$c((u_1,u_2,\cdots,u_k)+X)=(a_1(h_1u_1), a_2(h_2u_2), \cdots, a_k(h_ku_k)).$$
\\
\textbf{Claim 2}.  The above correspondence defines a group action of $G$ on $W$.

Assume that $c=a_1a_2\cdots a_k=a_1'a_2' \cdots a_k'$, where $a_i, a_i' \in A_i$. Write $h_i=a_{i+1}\cdots a_k, h_i'=a_{i+1}'\cdots a_k'$ for $1\leq i \leq k-1$ and $h_k=h_k'=1$. We also choose
$$(u_1,u_2,\cdots,u_k)+X=(u_1',u_2',\cdots,u_k')+X, u_i, u_i' \in U_i,$$
and we will show that $(a_1(h_1u_1)-a_1'(h_1'u_1'), \cdots, a_k(h_ku_k)-a_k'(h_k'u_k')) \in X$.
For each $1 \leq i \leq k$, it follows from Lemma~\ref{lem-elements} that $a_i^{-1}a_i' \in \Z(A_i)$ and $h_i^{-1}h_i' \in \C_G(A_iA_{i+1})$. Let $u \in U_i$. As $h^{-1}h_i'$ centralizes $A_i$ and $\pi_i^{-1}(u) \in A_i$, it follows that
$$(h_i^{-1}h_i')u=\pi_i((h_i^{-1}h_i')\pi_i^{-1}(u)(h_i^{-1}h_i')^{-1})=\pi_i(\pi_i^{-1}(u))=u,$$
that is, $h_i^{-1}h_i'$ acts trivially on $U$. Hence $h_i'u_i'=h_i(h_i^{-1}h_i')u_i'=h_iu_i'$. By hypothesis, $a_i^{-1}a_i' \in \Z(A_i) \subseteq \Ker(A_i~on~U_i)$. Hence
$$a_i'(h_i'u_i')=a_i'(h_iu_i')=a_i(a_i^{-1}a_i')(h_iu_i')=a_i(h_iu_i').$$
Since $(u_1-u_1',u_2-u_2',\cdots,u_k-u_k') \in X$, there exist $x_i \in A_i (1 \leq i \leq k)$ such that
$u_i-u_i'=\pi_i(x_i) (1\leq i \leq k)$ and $x_1x_2\cdots x_k=1$. By Lemma~\ref{lem-elements}(2), $x_i \in \Z(A_i)$ for each $i$ and hence ${}^{h_i}x_i \in \Z(A_i)$ as $H_i$ normalises $\Z(A_i)$.
Note that
$$
\begin{aligned}
a_i(h_iu_i)-a_i'(h_i'u_i')&=a_i(h_iu_i)-a_i(h_iu_i')=a_i(h_i(u_i-u_i'))\\
&=a_i(h_i\pi_i(x_i))=a_i\pi_i({}^{h_i}x_i)\\
&=\pi_i({}^{a_ih_i}x_i)=\pi_i({}^{h_i}x_i).
\end{aligned}
$$
It follows from Lemma~\ref{lem-elements}(3) that
$({}^{h_1}x_1)\cdot ({}^{h_2}x_2) \cdots ({}^{h_k}x_k)=1$, which implies that $(a_1(h_1u_1)-a_1'(h_1'u_1'), \cdots, a_k(h_ku_k)-a_k'(h_k'u_k')) \in X$. Hence the above correspondence is actually a map.

Next we show that this map defines an group action of $G$ on $W$. For any $b=b_1b_2\cdots b_k, c=c_1c_2\cdots c_k \in G$ and $(u_1,u_2, \cdots, u_k) +X \in W$,  where $b_i,c_i \in A_i, u_i \in U_i$,  we write $s_i=b_{i+1}b_{i+2} \cdots b_k, t_i=c_{i+1}c_{i+2} \cdots c_k$ for $1\leq i \leq k-1$ and $s_k=t_k=1$. Note that we easily see by induction,
$$
\begin{aligned}
bc&=(b_1s_1)(c_1t_1)=(b_1\cdot {}^{s_1}c_1)(s_1t_1)=(b_1\cdot {}^{s_1}c_1)(b_2s_2c_2t_2)\\
&=(b_1\cdot {}^{s_1}c_1)(b_2\cdot {}^{s_2}c_2)(s_2t_2)\\
&\cdots\\
&=(b_1\cdot {}^{s_1}c_1)(b_2\cdot {}^{s_2}c_2) \cdots (b_{i}\cdot {}^{s_{i}}c_{i})(s_it_i)\\
&\cdots\\
&=(b_1\cdot {}^{s_1}c_1)(b_2\cdot {}^{s_2}c_2) \cdots (b_{k-1}\cdot {}^{s_{k-1}}c_{k-1})(b_{k}\cdot {}^{s_{k}}c_{k}),
\end{aligned}
$$
and $b_i\cdot {}^{s_i}c_i \in A_i$ as $s_i \in H_i$ normalises $A_i$. We write
$(bc)((u_1,u_2, \cdots, u_k) +X)=(v_1,v_2, \cdots, v_k) +X$, where $v_i=(b_{i}\cdot {}^{s_{i}}c_{i})(s_it_iu_i)
=(b_is_ic_it_i)u_i.$ Then
$$
\begin{aligned}
b(c((u_1,u_2, \cdots, u_k) +X))&=b(((c_1t_1)u_1,(c_2t_2)u_2, \cdots, (c_ks_k)u_k) +X)\\
&=((b_1s_1c_1t_1)u_1,(b_2s_2c_2t_2)u_2, \cdots, (b_ks_kc_ks_k)u_k) +X\\
&=(v_1,v_2, \cdots, v_k) +X=(bc)((u_1,u_2, \cdots, u_k) +X).
\end{aligned}
$$
Consequently, we have an action of the group $G$ on the $A$-module $W$.
\\
\\
Consider the correspondence $\pi: G \rightarrow W$ given by
$$\pi(g)=(\pi_1(a_1), \pi_2(a_2), \cdots, \pi_k(a_k))+X,$$
where $g=a_1a_2\cdots a_k$ and $a_i \in A_i (1\leq i \leq k)$.
\\
\\
\textbf{Claim 3}.  The correspondence $\pi$ is a bijective map from $G$ to $W$.

We first show that $\pi$ is a map. Let $g=a_1a_2\cdots a_k=b_1b_2\cdots b_k$, where $a_i, b_i \in A_i, 1\leq i \leq k$. By Lemma~\ref{lem-elements}(2) and Hypothesis (3), $a_ib^{-1}_i \in \Z(A_i) \subseteq \Ker(A_i~on~U_i)$. It follows from Lemma~\ref{lem-1-cocycle} that $\pi_i(a_i)=\pi_i((a_ib_i^{-1})b_i)=\pi_i(a_ib_i^{-1})+\pi_i(b_i)$, that is, $\pi_i(a_i)-\pi_i(b_i)=\pi_i(a_ib_i^{-1})$. Hence
$$
\begin{aligned}
&(\pi_1(a_1), \pi_2(a_2), \cdots, \pi_k(a_k))-(\pi_1(b_1), \pi_2(b_2), \cdots, \pi_k(b_k))\\
=&(\pi_1(a_1)-\pi_1(b_1), \pi_2(a_2)-\pi_2(b_2), \cdots, \pi_k(a_k)-\pi_k(b_k))\\
=&(\pi_1(a_1b_1^{-1}), \pi_2(a_2b_2^{-1}), \cdots, \pi_k(a_kb_k^{-1})) \in X;
\end{aligned}
$$
as $(a_1b_1^{-1}) \cdot (a_2b_2^{-1}) \cdots (a_kb_k^{-1})=1$ by Lemma~\ref{lem-elements}(1). Thus $\pi$ is a map. For each $(\pi_1(a_1), \pi_2(a_2), \cdots, \pi_k(a_k))+X \in W$, where $a_i \in A_i$, we take $g=a_1a_2\cdots a_k \in G$. Then $\pi(g)=(\pi_1(a_1), \pi_2(a_2), \cdots, \pi_k(a_k))+X$. Hence $\pi$ is surjective. Write
$$\XX_k=\{(a_1,a_2,\cdots,a_k) \mid a_1a_2\cdots a_k=1, a_i \in A_i, \forall 1\leq i \leq k\}.$$
Since $\pi_i$ is bijective for each $i$, it follows that $|\XX_k|=|X|$. Hence, by Lemma~\ref{lem-elements}(4), we have
$$|G|=|A_1A_2 \cdots A_k|=|A_1||A_2| \cdots |A_k|/|\XX_k|=|U_1||U_2| \cdots |U_k|/|X|=|W|.$$
Thus $\pi$ is a bijective map from $G$ to $W$.
\\
\\
\textbf{Claim 4}.  The map $\pi$ is a bijective $1$-cocycle from $G$ to $W$ with respect to the action of $G$ on $W$.

Let $b=b_1b_2\cdots b_k, c=c_1c_2\cdots c_k \in G$,  where $b_i,c_i \in A_i$ and write $s_i=b_{i+1}b_{i+2} \cdots b_k, t_i=c_{i+1}c_{i+2} \cdots c_k$ for $1\leq i \leq k-1$ and $s_k=t_k=1$.
Note that $bc=(b_1\cdot {}^{s_1}c_1)(b_2\cdot {}^{s_2}c_2) \cdots (b_{k-1}\cdot {}^{s_{k-1}}c_{k-1})(b_{k}\cdot {}^{s_{k}}c_{k})$ and $b_{i}\cdot {}^{s_{i}}c_{i} \in A_i$. Then
$$
\begin{aligned}
\pi(bc)&=\pi((b_1\cdot {}^{s_1}c_1)(b_2\cdot {}^{s_2}c_2) \cdots (b_{k-1}\cdot {}^{s_{k-1}}c_{k-1})(b_{k}\cdot {}^{s_{k}}c_{k}))\\
&=(\pi_1(b_1\cdot {}^{s_1}c_1), \pi_2(b_2\cdot {}^{s_2}c_2), \cdots, \pi_k(b_k\cdot {}^{s_k}c_k))+X\\
&=(\pi_1(b_1)+b_1\pi_1({}^{s_1}c_1), \pi_2(b_2)+b_2\pi_2({}^{s_2}c_2), \cdots, \pi_k(b_k)+b_k\pi_k({}^{s_k}c_k))+X\\
&=\pi(b)+(b_1s_1\pi_1(c_1), b_2s_2\pi_2(c_2), \cdots, b_ks_k\pi_k(c_k))+X\\
&=\pi(b)+b\pi(c).
\end{aligned}
$$
Hence $(W,\pi)$ is an IYB-structure on $G$. We bring the proof to an end proving that
$(W,\pi)$ is $A$-equivariant. For each $a \in A$ and $g=g_1\cdots g_k \in G$, where $g_i \in A_i, 1\leq i \leq k$, we have that
$$
\begin{aligned}
a\pi(g)&=a((\pi_1(g_1), \cdots, \pi_k(g_k))+X)\\
&=(a\pi_1(g_1), \cdots, a\pi_k(g_k))+X\\
&=(\pi_1({}^ag_1), \cdots, \pi_k({}^ag_k))+X\\
&=\pi({}^ag_1\cdots {}^ag_k)=\pi({}^a(g_1\cdots g_k))=\pi({}^ag),
\end{aligned}
$$
as desired. The proof of the theorem is now complete.
\end{proof}



%Let a group $G$ act on an elementary abelian $p$-group $V$ for some prime $p$ via automorphisms. An element $a \in G$ acts \emph{quadratically} on $V$ if $[V,a,a]=1$. The action of $G$ on $V$ is called \emph{$p$-stable} if
%$a\C_G(V) \in \Oo_p(G/\C_G(V))$ for each element $a \in G$ such that $[V,a,a]=1$.

%\begin{lemma}
%Let a group $G$ act faithfully on an elementary abelian $p$-group $V$ for some prime $p \neq 2$. Suppose that the action of $G$ on $V$ is not $p$-stable. Then
%\begin{itemize}
%\item[(a)] The Sylow $2$-subgroups of $G$ are non-Abelian.
%\item[(b)] If in addition $G$ is $p$-soluble, then $p=3$ and there exists a section of $G$ isomorphic to $\SL_2(3)$.
%\end{itemize}
%\end{lemma}

As a corollary of Theorem~\ref{thm-main}, the following extension of Theorem~\ref{thm-Q-8} holds.
\begin{corollary}\label{cor-semidirect}
Let $G$ be a finite soluble group and let $N$ be a normal Hall subgroup of $G$ such that $G/N$ is an IYB-group. Suppose that all Sylow subgroups of $G$ have nilpotency class at most two and the nilpotent residual $N^{\mathcal{N}}$ of $N$ is $\Q_8$-free.
Then $G$ is an IYB-group.
\end{corollary}

%The proof of Theorem~\ref{thm-main} is given in Section~2, where we also collect some results on IYB-structures. In Section~3, we introduce system normalisers in the context of finite soluble groups. Finally, Section~4 presents applications of our framework, including the proofs of Theorem~\ref{thm-Q-8} and Corollary~\ref{cor-semidirect}.

\begin{proof}[\textbf{\emph{Proof of Corollary~\ref{cor-semidirect}}}]
By Schur-Zassenhaus Theorem~, $N$ has a complement $A$ in $G$. Considering the coprime action of $A$ on $N$ via conjugation, it follows from Lemma~\ref{lem-nilpotency-class-2} , there exist
some $A$-invariant nilpotent subgroups $A_1,\cdots,A_k$ of $G$ such that
\begin{itemize}
\item[(1)] $N=A_1A_2\cdots A_k$;
\item[(2)] $A_i$ is normalised by $A_j$ for $1\leq i <j\leq k$;
\item[(3)] $A_1A_2\cdots A_i \cap A_{i+1}\cdots A_k \leq \C_G(A_iA_{i+1})$ for each $1 \leq i \leq k-1$;
\item[(4)] $A_i$ has nilpotency class at most two with $|A_i'|$ odd for each $1\leq i \leq k$.
\end{itemize}
Write $H_i=A_{i+1}A_{i+2}\cdots A_k$ for $i=1,2,\cdots, k-1$ and Clearly $A_i$ is normalised by $H_i$ by Condition $(2)$. By (4), it follows from Lemma~\ref{lem-odd-order-commutator} that $A_i$ has a fully equivariant IYB-structure $(U_i,\pi_i)$ with $\Z(A_i) \subseteq \Ker(A_i~on~U_i)$.
Hence $(U_i,\pi_i)$ is $AH_i$-invariant.
Applying Theorem~\ref{thm-main}, $N$ has an $A$-equivariant IYB-structure. By hypothesis, $A \cong G/N$ is an IYB-group. It follows from Lemma~\ref{lem-equivariant-IYB} that $G$ is an IYB-group.
\end{proof}

